% Options for packages loaded elsewhere
\PassOptionsToPackage{unicode}{hyperref}
\PassOptionsToPackage{hyphens}{url}
%
\documentclass[
  11pt,
]{article}
\usepackage{amsmath,amssymb}
\usepackage{iftex}
\ifPDFTeX
  \usepackage[T1]{fontenc}
  \usepackage[utf8]{inputenc}
  \usepackage{textcomp} % provide euro and other symbols
\else % if luatex or xetex
  \usepackage{unicode-math} % this also loads fontspec
  \defaultfontfeatures{Scale=MatchLowercase}
  \defaultfontfeatures[\rmfamily]{Ligatures=TeX,Scale=1}
\fi
\usepackage{lmodern}
\ifPDFTeX\else
  % xetex/luatex font selection
\fi
% Use upquote if available, for straight quotes in verbatim environments
\IfFileExists{upquote.sty}{\usepackage{upquote}}{}
\IfFileExists{microtype.sty}{% use microtype if available
  \usepackage[]{microtype}
  \UseMicrotypeSet[protrusion]{basicmath} % disable protrusion for tt fonts
}{}
\makeatletter
\@ifundefined{KOMAClassName}{% if non-KOMA class
  \IfFileExists{parskip.sty}{%
    \usepackage{parskip}
  }{% else
    \setlength{\parindent}{0pt}
    \setlength{\parskip}{6pt plus 2pt minus 1pt}}
}{% if KOMA class
  \KOMAoptions{parskip=half}}
\makeatother
\usepackage{xcolor}
\usepackage[a4paper,margin=2.5cm]{geometry}
\usepackage{longtable,booktabs,array}
\usepackage{calc} % for calculating minipage widths
% Correct order of tables after \paragraph or \subparagraph
\usepackage{etoolbox}
\makeatletter
\patchcmd\longtable{\par}{\if@noskipsec\mbox{}\fi\par}{}{}
\makeatother
% Allow footnotes in longtable head/foot
\IfFileExists{footnotehyper.sty}{\usepackage{footnotehyper}}{\usepackage{footnote}}
\makesavenoteenv{longtable}
\usepackage{graphicx}
\makeatletter
\def\maxwidth{\ifdim\Gin@nat@width>\linewidth\linewidth\else\Gin@nat@width\fi}
\def\maxheight{\ifdim\Gin@nat@height>\textheight\textheight\else\Gin@nat@height\fi}
\makeatother
% Scale images if necessary, so that they will not overflow the page
% margins by default, and it is still possible to overwrite the defaults
% using explicit options in \includegraphics[width, height, ...]{}
\setkeys{Gin}{width=\maxwidth,height=\maxheight,keepaspectratio}
% Set default figure placement to htbp
\makeatletter
\def\fps@figure{htbp}
\makeatother
\setlength{\emergencystretch}{3em} % prevent overfull lines
\providecommand{\tightlist}{%
  \setlength{\itemsep}{0pt}\setlength{\parskip}{0pt}}
\setcounter{secnumdepth}{-\maxdimen} % remove section numbering
\ifLuaTeX
\usepackage[bidi=basic]{babel}
\else
\usepackage[bidi=default]{babel}
\fi
\babelprovide[main,import]{norsk}
% get rid of language-specific shorthands (see #6817):
\let\LanguageShortHands\languageshorthands
\def\languageshorthands#1{}
\usepackage{pifont}
\usepackage{seqsplit}
\usepackage{newunicodechar}
\usepackage{caption}
\usepackage{pdflscape}
\usepackage{microtype}
\usepackage{etoolbox}
\usepackage{array}
\usepackage[export]{adjustbox}

% Global takhøyde for figurer. Pandoc-spesifisert width=NN% beholdes,
% men keepaspectratio garanterer proporsjonal skalering, og max
% totalheight stopper figuren i å renne over sidekanten dersom pandoc
% emitterer height=\textheight (som er default i pandoc 3.x).
\setkeys{Gin}{keepaspectratio,max width=\linewidth,max totalheight=0.8\textheight}

% Slakkere LaTeX-float-parametre slik at figurer lettere finner en
% gyldig plassering når siden allerede er full av tekst og tabeller.
% Default-verdiene fører ellers til Overfull \vbox når en moderat
% høy figur (≈9 cm) presses inn nederst på en side med liten ledig
% plass i stedet for å flyte til neste side.
\renewcommand{\topfraction}{0.9}
\renewcommand{\bottomfraction}{0.9}
\renewcommand{\textfraction}{0.08}
\renewcommand{\floatpagefraction}{0.7}

\setlength{\tabcolsep}{4pt}
\renewcommand{\arraystretch}{1.1}
\AtBeginEnvironment{longtable}{\small}

% Caption-styling: ingen auto-prefiks (vi har manuell "Figur X.Y" /
% "Tabell X.Y" i caption-teksten), liten italic skrift sentrert under
% figur og tabell.
\captionsetup[figure]{labelformat=empty, font={small,it}, justification=centering}
\captionsetup[longtable]{labelformat=empty, font={small,it}, justification=centering}

% Tillat linjebryting i lange ubrudte strenger (klassenavn, parametre osv.)
% slik at de ikke renner over kolonnegrensen i smale tabellceller. Visuelt
% nøytralt når strengen får plass uten bryting.
\newcommand{\splitcode}[1]{{\ttfamily\seqsplit{#1}}}
\newcommand{\splitword}[1]{\seqsplit{#1}}

% Default-fonten Latin Modern har ikke disse Unicode-symbolene; mapp dem
% til pifont-glyfer / LaTeX-ekvivalenter.
\newunicodechar{☐}{\ding{110}}
\newunicodechar{☒}{\ding{55}}
\newunicodechar{−}{\ensuremath{-}}
\ifLuaTeX
  \usepackage{selnolig}  % disable illegal ligatures
\fi
\IfFileExists{bookmark.sty}{\usepackage{bookmark}}{\usepackage{hyperref}}
\IfFileExists{xurl.sty}{\usepackage{xurl}}{} % add URL line breaks if available
\urlstyle{same}
\hypersetup{
  pdflang={nb-NO},
  hidelinks,
  pdfcreator={LaTeX via pandoc}}

\author{}
\date{}

\begin{document}

\hypertarget{salgsprognoser-for-en-simulert-dagligvarekjede-en-sammenligning-av-lineuxe6r-regresjon-og-random-forest-regressor}{%
\section{\texorpdfstring{Salgsprognoser for en simulert dagligvarekjede: en sammenligning av lineær regresjon og Random
Forest Regressor
}{Salgsprognoser for en simulert dagligvarekjede: en sammenligning av lineær regresjon og Random Forest Regressor }}\label{salgsprognoser-for-en-simulert-dagligvarekjede-en-sammenligning-av-lineuxe6r-regresjon-og-random-forest-regressor}}

\hypertarget{del-av-log650-forskningsprosjekt-logistikk-og-kunstig-intelligens}{%
\subsection{\texorpdfstring{Del av LOG650 Forskningsprosjekt: Logistikk og Kunstig Intelligens
}{Del av LOG650 Forskningsprosjekt: Logistikk og Kunstig Intelligens }}\label{del-av-log650-forskningsprosjekt-logistikk-og-kunstig-intelligens}}

\begin{center}

\includegraphics{extracted_images/image_0.png}
\end{center}

Forfatter(e): Erik Brendehaug, Joseph James, Pål Rånes, Marthe Slåtta Bjerke

Totalt antall sider inkludert forsiden: TODO (fylles inn ved endelig låsing)

Molde, Innleveringsdato: TODO (fylles inn nærmere innlevering)

\begin{center}

\includegraphics{extracted_images/image_1.png}
\end{center}

\newpage

\hypertarget{obligatorisk-egenerkluxe6ringgruppeerkluxe6ring}{%
\subsection{\texorpdfstring{Obligatorisk egenerklæring/gruppeerklæring
}{Obligatorisk egenerklæring/gruppeerklæring }}\label{obligatorisk-egenerkluxe6ringgruppeerkluxe6ring}}

Den enkelte student er selv ansvarlig for å sette seg inn i hva som er lovlige hjelpemidler, retningslinjer for bruk av
disse og regler om kildebruk.

Du/dere fyller ut erklæringen ved å klikke i ruten til høyre for den enkelte del 1--6:

\begin{longtable}[]{@{}
  >{\raggedright\arraybackslash}p{(\columnwidth - 4\tabcolsep) * \real{0.3333}}
  >{\raggedright\arraybackslash}p{(\columnwidth - 4\tabcolsep) * \real{0.3333}}
  >{\raggedright\arraybackslash}p{(\columnwidth - 4\tabcolsep) * \real{0.3333}}@{}}
\toprule\noalign{}
\begin{minipage}[b]{\linewidth}\raggedright
Nr.
\end{minipage} & \begin{minipage}[b]{\linewidth}\raggedright
Erklæring
\end{minipage} & \begin{minipage}[b]{\linewidth}\raggedright
Avkryssing
\end{minipage} \\
\midrule\noalign{}
\endhead
\bottomrule\noalign{}
\endlastfoot
1 & Jeg/vi erklærer at min/vår besvarelse er mitt/vårt eget arbeid. & ☒ \\
2 & Jeg/vi erklærer videre at denne besvarelsen:- ikke har vært brukt til annen eksamen- ikke refererer til andres
arbeid uten at det er oppgitt- ikke refererer til eget tidligere arbeid uten at det er oppgitt- har alle referanser
oppgitt- ikke er kopi eller duplikat & ☒ \\
3 & Brudd på ovennevnte er fusk og kan medføre annullering. & ☒ \\
4 & Oppgaven kan bli \splitword{plagiatkontrollert} i URKUND. & ☒ \\
5 & Høgskolen vil behandle mistanke om fusk etter \splitword{retningslinjene.} & ☒ \\
6 & Jeg/vi har satt oss inn i reglene for kildebruk. & ☒ \\
\end{longtable}

\hypertarget{section}{%
\subsection{}\label{section}}

\hypertarget{personvern}{%
\subsection{\texorpdfstring{Personvern }{Personvern }}\label{personvern}}

\hypertarget{personopplysningsloven}{%
\subsubsection{\texorpdfstring{Personopplysningsloven }{Personopplysningsloven }}\label{personopplysningsloven}}

\begin{longtable}[]{@{}
  >{\raggedright\arraybackslash}p{(\columnwidth - 4\tabcolsep) * \real{0.3333}}
  >{\raggedright\arraybackslash}p{(\columnwidth - 4\tabcolsep) * \real{0.3333}}
  >{\raggedright\arraybackslash}p{(\columnwidth - 4\tabcolsep) * \real{0.3333}}@{}}
\toprule\noalign{}
\begin{minipage}[b]{\linewidth}\raggedright
Spørsmål
\end{minipage} & \begin{minipage}[b]{\linewidth}\raggedright
Ja
\end{minipage} & \begin{minipage}[b]{\linewidth}\raggedright
Nei
\end{minipage} \\
\midrule\noalign{}
\endhead
\bottomrule\noalign{}
\endlastfoot
Har oppgaven vært vurdert av NSD? & ☐ & ☒ \\
\end{longtable}

\begin{itemize}
\item
  Hvis ja:\\
  Referansenummer:
\item
  Hvis nei:\\
  Jeg/vi erklærer at oppgaven ikke omfattes av Personopplysningsloven ☒
\end{itemize}

\begin{center}\rule{0.5\linewidth}{0.5pt}\end{center}

\hypertarget{helseforskningsloven}{%
\subsubsection{\texorpdfstring{Helseforskningsloven }{Helseforskningsloven }}\label{helseforskningsloven}}

\begin{longtable}[]{@{}
  >{\raggedright\arraybackslash}p{(\columnwidth - 4\tabcolsep) * \real{0.3333}}
  >{\raggedright\arraybackslash}p{(\columnwidth - 4\tabcolsep) * \real{0.3333}}
  >{\raggedright\arraybackslash}p{(\columnwidth - 4\tabcolsep) * \real{0.3333}}@{}}
\toprule\noalign{}
\begin{minipage}[b]{\linewidth}\raggedright
Spørsmål
\end{minipage} & \begin{minipage}[b]{\linewidth}\raggedright
Ja
\end{minipage} & \begin{minipage}[b]{\linewidth}\raggedright
Nei
\end{minipage} \\
\midrule\noalign{}
\endhead
\bottomrule\noalign{}
\endlastfoot
Har oppgaven vært til behandling hos REK? & ☐ & ☒ \\
\end{longtable}

\begin{itemize}
\tightlist
\item
  Hvis ja:\\
  Referansenummer:
\end{itemize}

\begin{center}\rule{0.5\linewidth}{0.5pt}\end{center}

\hypertarget{publiseringsavtale}{%
\subsection{\texorpdfstring{Publiseringsavtale }{Publiseringsavtale }}\label{publiseringsavtale}}

\begin{longtable}[]{@{}
  >{\raggedright\arraybackslash}p{(\columnwidth - 2\tabcolsep) * \real{0.5000}}
  >{\raggedright\arraybackslash}p{(\columnwidth - 2\tabcolsep) * \real{0.5000}}@{}}
\toprule\noalign{}
\begin{minipage}[b]{\linewidth}\raggedright
Felt
\end{minipage} & \begin{minipage}[b]{\linewidth}\raggedright
Verdi
\end{minipage} \\
\midrule\noalign{}
\endhead
\bottomrule\noalign{}
\endlastfoot
Studiepoeng & 15 \\
Veileder & Bård Inge Peterson (BIP) \\
\end{longtable}

\hypertarget{fullmakt-til-elektronisk-publisering}{%
\subsubsection{\texorpdfstring{Fullmakt til elektronisk publisering
}{Fullmakt til elektronisk publisering }}\label{fullmakt-til-elektronisk-publisering}}

\begin{longtable}[]{@{}
  >{\raggedright\arraybackslash}p{(\columnwidth - 4\tabcolsep) * \real{0.3333}}
  >{\raggedright\arraybackslash}p{(\columnwidth - 4\tabcolsep) * \real{0.3333}}
  >{\raggedright\arraybackslash}p{(\columnwidth - 4\tabcolsep) * \real{0.3333}}@{}}
\toprule\noalign{}
\begin{minipage}[b]{\linewidth}\raggedright
Spørsmål
\end{minipage} & \begin{minipage}[b]{\linewidth}\raggedright
Ja
\end{minipage} & \begin{minipage}[b]{\linewidth}\raggedright
Nei
\end{minipage} \\
\midrule\noalign{}
\endhead
\bottomrule\noalign{}
\endlastfoot
Gjøre oppgaven tilgjengelig for elektronisk publisering & ☒ & ☐ \\
\end{longtable}

\begin{longtable}[]{@{}
  >{\raggedright\arraybackslash}p{(\columnwidth - 4\tabcolsep) * \real{0.3333}}
  >{\raggedright\arraybackslash}p{(\columnwidth - 4\tabcolsep) * \real{0.3333}}
  >{\raggedright\arraybackslash}p{(\columnwidth - 4\tabcolsep) * \real{0.3333}}@{}}
\toprule\noalign{}
\begin{minipage}[b]{\linewidth}\raggedright
Spørsmål
\end{minipage} & \begin{minipage}[b]{\linewidth}\raggedright
Ja
\end{minipage} & \begin{minipage}[b]{\linewidth}\raggedright
Nei
\end{minipage} \\
\midrule\noalign{}
\endhead
\bottomrule\noalign{}
\endlastfoot
Er oppgaven båndlagt \splitword{(konfidensiell)?} & ☐ & ☒ \\
\end{longtable}

\begin{itemize}
\tightlist
\item
  Hvis ja:
\end{itemize}

\begin{longtable}[]{@{}
  >{\raggedright\arraybackslash}p{(\columnwidth - 4\tabcolsep) * \real{0.3333}}
  >{\raggedright\arraybackslash}p{(\columnwidth - 4\tabcolsep) * \real{0.3333}}
  >{\raggedright\arraybackslash}p{(\columnwidth - 4\tabcolsep) * \real{0.3333}}@{}}
\toprule\noalign{}
\begin{minipage}[b]{\linewidth}\raggedright
Spørsmål
\end{minipage} & \begin{minipage}[b]{\linewidth}\raggedright
Ja
\end{minipage} & \begin{minipage}[b]{\linewidth}\raggedright
Nei
\end{minipage} \\
\midrule\noalign{}
\endhead
\bottomrule\noalign{}
\endlastfoot
Kan publiseres etter \splitword{båndleggingsperiode} & ☐ & ☐ \\
\end{longtable}

Dato: TODO (fylles inn nærmere innlevering)

\begin{center}\rule{0.5\linewidth}{0.5pt}\end{center}

\hypertarget{antall-ord}{%
\subsection{\texorpdfstring{Antall ord }{Antall ord }}\label{antall-ord}}

Rapportens hovedtekst (kap. 1 Innledning til og med kap. 10 Konklusjon) er på omtrent 5 300 ord. Sammendrag, abstract,
forside, litteraturliste og vedlegg er holdt utenfor tellingen.

\hypertarget{forfattererkluxe6ring}{%
\subsection{\texorpdfstring{Forfattererklæring }{Forfattererklæring }}\label{forfattererkluxe6ring}}

Gruppen har i fellesskap skrevet og godkjent denne rapporten.

\begin{center}\rule{0.5\linewidth}{0.5pt}\end{center}

\hypertarget{sammendrag}{%
\subsection{\texorpdfstring{Sammendrag }{Sammendrag }}\label{sammendrag}}

Dagligvarehandelens driftsrisiko er direkte knyttet til kvaliteten på etterspørselsprognosene: for høye bestillinger gir
svinn og bundet kapital, mens for lave gir tomme hyller og tapte inntekter. Denne rapporten undersøker hvordan multippel
lineær regresjon og Random Forest Regressor kan brukes til å predikere daglig salg for 2025 i en simulert
dagligvarekjede, og hvilke faktorer som påvirker salget mest.

Analysen bygger på et datasett med 9 994 daglige salgstransaksjoner fra perioden 2022--2025. Treningsperioden 2022--2024
dekker 6 682 rader og testperioden 2025 dekker 3 312 rader, med tidsbasert splitt for å unngå datalekkasje.
Modellmatrisen består av 67 features som inkluderer rabatt, kategoriske dimensjoner som region, produktkategori og
subkategori, samt sju utledede kalendervariabler. Tre modellspor evalueres på 2025: en benchmark multippel lineær
regresjon, en Random Forest-baseline med standardparametere, og en tuned Random Forest der hyperparametere er valgt
gjennom et rutenettsøk med 2022--2023 som søketrening og 2024 som valideringsår. Modellene sammenlignes med RMSE som
primær metrikk og MAPE som sekundær metrikk, både samlet, per måned og per segment fordelt på kvartal, rabattband,
region og salgsnivå.

Tuned Random Forest er den samlet beste modellen med RMSE 578,26 og MAPE 43,97 \%, og vinner RMSE i elleve av tolv
måneder og tretten av fjorten segmenter. Gapet til benchmark lineær er marginalt på totalnivå, og benchmark lineær er
best i segmentet for høyt salgsnivå. De mest påvirkningsrike prediktorene er rabatt (11,48 \%) og kalendervariablene,
der kalendervariablene alene utgjør om lag 42 \% av importance i topp 10. Rapporten anbefaler tuned Random Forest som
standardprognose for innkjøp, lager og aggregert ressursplanlegging, benchmark lineær som forklaringsstøtte, og baseline
Random Forest som kontroll mot prosentfeil i høyrabattperioder. Funnene er prediktive og ikke kausale, og gjelder
innenfor det simulerte caset.

\hypertarget{abstract}{%
\subsection{\texorpdfstring{Abstract }{Abstract }}\label{abstract}}

Operational risk in grocery retail is directly linked to forecast quality: over-ordering leads to waste and tied-up
capital, while under-ordering causes stock-outs and lost revenue. This report examines how multiple linear regression
and Random Forest Regressor can be used to predict daily sales for 2025 in a simulated grocery chain, and which factors
influence sales the most.

The analysis uses 9,994 daily sales transactions from 2022--2025. Training data covers 2022--2024 (6,682 rows) and test
data covers 2025 (3,312 rows), with a time-based split to prevent data leakage. The model matrix contains 67 features,
including discount, categorical dimensions such as region, product category and sub-category, and seven derived calendar
variables. Three tracks are evaluated on 2025: a benchmark multiple linear regression, a Random Forest baseline with
default parameters, and a tuned Random Forest whose hyperparameters are selected via grid search on 2022--2023 with 2024
as the validation year. RMSE is the primary metric and MAPE the secondary, both overall, by month and by segment across
quarter, discount band, region and sales level.

The tuned Random Forest is the overall best model with RMSE 578.26 and MAPE 43.97 \%, winning RMSE in eleven of twelve
months and thirteen of fourteen segments. The gap to the benchmark linear model is marginal overall, and the benchmark
linear model is the strongest in the high-sales segment. The most influential predictors are discount (11.48 \%) and the
calendar variables, where the calendar variables alone account for around 42 \% of importance in the top ten. The report
recommends the tuned Random Forest as the standard forecast for purchasing, inventory and aggregate resource planning,
the benchmark linear model as explanatory support, and the Random Forest baseline as a check against percentage error
during high-discount periods. The findings are predictive rather than causal and apply within the simulated case.

\begin{center}\rule{0.5\linewidth}{0.5pt}\end{center}

\hypertarget{innhold}{%
\subsection{\texorpdfstring{Innhold }{Innhold }}\label{innhold}}

\begin{itemize}
\tightlist
\item
  \protect\hyperlink{1-innledning}{1 Innledning}

  \begin{itemize}
  \tightlist
  \item
    \protect\hyperlink{11-problemstilling}{1.1 Problemstilling}
  \item
    \protect\hyperlink{12-delproblemer}{1.2 Delproblemer}
  \item
    \protect\hyperlink{13-avgrensinger}{1.3 Avgrensinger}
  \item
    \protect\hyperlink{14-antagelser}{1.4 Antagelser}
  \end{itemize}
\item
  \protect\hyperlink{2-litteratur}{2 Litteratur}
\item
  \protect\hyperlink{3-teori}{3 Teori}

  \begin{itemize}
  \tightlist
  \item
    \protect\hyperlink{31-multippel-lineuxe6r-regresjon}{3.1 Multippel lineær regresjon}
  \item
    \protect\hyperlink{32-random-forest-regressor}{3.2 Random Forest Regressor}
  \item
    \protect\hyperlink{33-evalueringsmetrikker}{3.3 Evalueringsmetrikker}
  \item
    \protect\hyperlink{34-feature-engineering-og-dataoppsett}{3.4 Feature engineering og dataoppsett}
  \end{itemize}
\item
  \protect\hyperlink{4-casebeskrivelse}{4 Casebeskrivelse}

  \begin{itemize}
  \tightlist
  \item
    \protect\hyperlink{41-dagligvare-og-beslutningssituasjonen}{4.1 Dagligvare og beslutningssituasjonen}
  \item
    \protect\hyperlink{42-historisk-salgsutvikling}{4.2 Historisk salgsutvikling}
  \item
    \protect\hyperlink{43-sesongmuxf8nster-i-salget}{4.3 Sesongmønster i salget}
  \item
    \protect\hyperlink{44-utfordringer-duxe5rlige-prognoser-medfuxf8rer-i-bedriften}{4.4 Utfordringer dårlige prognoser
    medfører i bedriften}
  \end{itemize}
\item
  \protect\hyperlink{5-metode-og-data}{5 Metode og data}

  \begin{itemize}
  \tightlist
  \item
    \protect\hyperlink{51-metode}{5.1 Metode}
  \item
    \protect\hyperlink{52-data}{5.2 Data}
  \end{itemize}
\item
  \protect\hyperlink{6-modellering}{6 Modellering}
\item
  \protect\hyperlink{7-analyse}{7 Analyse}
\item
  \protect\hyperlink{8-resultat}{8 Resultat}
\item
  \protect\hyperlink{9-diskusjon}{9 Diskusjon}

  \begin{itemize}
  \tightlist
  \item
    \protect\hyperlink{91-tolkning-av-hovedfunn-mot-problemstillingen}{9.1 Tolkning av hovedfunn mot problemstillingen}
  \item
    \protect\hyperlink{92-variablenes-puxe5virkning}{9.2 Variablenes påvirkning}
  \item
    \protect\hyperlink{93-praktisk-nytte-for-dagligvare}{9.3 Praktisk nytte for Dagligvare}
  \item
    \protect\hyperlink{94-metodiske-begrensninger}{9.4 Metodiske begrensninger}
  \item
    \protect\hyperlink{95-videre-arbeid}{9.5 Videre arbeid}
  \end{itemize}
\item
  \protect\hyperlink{10-konklusjon}{10 Konklusjon}
\item
  \protect\hyperlink{11-bibliografi}{11 Bibliografi}
\item
  \protect\hyperlink{12-vedlegg}{12 Vedlegg}
\end{itemize}

\begin{center}\rule{0.5\linewidth}{0.5pt}\end{center}

\hypertarget{innledning}{%
\subsection{1 Innledning}\label{innledning}}

Dagligvarehandel er preget av høy variabilitet i etterspørsel. Salget svinger med sesong, kampanjer, rabatter og
regionale forhold, og konsekvensene av dårlige prognoser er direkte synlige i driften: for høye bestillinger fører til
svinn og bundet kapital, mens for lave bestillinger gir utsolgte varer og tapte inntekter. Behovet for mer presise
etterspørselsprognoser er derfor sentralt når virksomheter skal fatte beslutninger om innkjøp, lager, kampanjer og
ressursplanlegging.

Maskinlæringsbaserte metoder som Random Forest Regressor gir nye muligheter for å fange opp ikke-lineære mønstre i
historiske salgsdata sammenlignet med tradisjonelle lineære tilnærminger. Dette prosjektet analyserer et simulert
dagligvarecase for perioden 2022--2025 med mål om å utvikle og evaluere prognosemodeller for salg i 2025, og å
identifisere hvilke forklaringsvariabler som bærer mest prediksjonsverdi.

Med dette som utgangspunkt formuleres problemstillingen som styrer metodevalg, analyse og tolkning videre i rapporten.

\hypertarget{problemstilling}{%
\subsubsection{1.1 Problemstilling}\label{problemstilling}}

Hvordan kan multippel lineær regresjon og Random Forest Regressor brukes til å forutsi salg for 2025 for en simulert
dagligvarekjede, og hvilke faktorer påvirker salget mest?

\hypertarget{delproblemer}{%
\subsubsection{1.2 Delproblemer}\label{delproblemer}}

Problemstillingen dekomponeres i to delproblemer som strukturerer analyse, diskusjon og konklusjon:

\begin{enumerate}
\def\labelenumi{\arabic{enumi}.}
\tightlist
\item
  \textbf{Modellvalg og prognoseytelse.} Hvordan sammenlignes de to modellklassene på 2025-data, målt med RMSE som
  primær og MAPE som sekundær metrikk, både samlet, månedlig og på tvers av segmenter som kvartal, rabattband, region og
  salgsnivå?
\item
  \textbf{Variabelanalyse.} Hvilke forklaringsvariabler har størst prediksjonsverdi, og hvordan samsvarer rangeringen
  fra Random Forest Regressor med koeffisientene i den lineære benchmarken?
\end{enumerate}

\hypertarget{avgrensinger}{%
\subsubsection{1.3 Avgrensinger}\label{avgrensinger}}

Analysen er avgrenset på følgende punkter:

\begin{enumerate}
\def\labelenumi{\arabic{enumi}.}
\tightlist
\item
  \textbf{Prediksjon, ikke kausal analyse.} Prosjektet undersøker hvilke variabler som predikerer salget best, ikke
  hvorfor salget endrer seg. Kausale slutninger forutsetter eksperimentelt design og ligger utenfor prosjektets rammer.
\item
  \textbf{Ingen lageroptimalisering.} Rapporten dokumenterer ikke lageroptimalisering. Det krever kostnads- og
  volumparametere som ikke inngår i datasettet.
\item
  \textbf{Ingen makroøkonomiske faktorer.} Eksterne variabler som inflasjon, rente og konjunkturer er ikke inkludert.
  Datasettet gir ikke grunnlag for dette, og inkludering av slike variabler ville kreve egne datakilder.
\item
  \textbf{Én simulert virksomhet og ett datasett.} Analysen begrenses til det tilgjengelige datasettet. Generalisering
  til andre virksomheter eller bransjer forutsetter eget datagrunnlag og ny validering.
\item
  \textbf{To modelltyper.} Modellvalget er begrenset til multippel lineær regresjon og Random Forest Regressor, i tråd
  med prosjektets faglige rammer og tilgjengelig treningsgrunnlag.
\end{enumerate}

\hypertarget{antagelser}{%
\subsubsection{1.4 Antagelser}\label{antagelser}}

Analysen bygger på følgende eksplisitte antagelser:

\begin{enumerate}
\def\labelenumi{\arabic{enumi}.}
\tightlist
\item
  \textbf{Representativitet.} Datasettet antas representativt for en realistisk dagligvarekjede med tanke på
  sesongmønstre, rabattbruk, regionfordeling og produktkategorier. \emph{Konsekvens:} Funnene er gyldige innenfor dette
  caset, men ikke nødvendigvis overførbare til andre virksomheter eller reelle kjeder.
\item
  \textbf{Tidsmessig relevans.} Historiske mønstre fra 2022--2024 antas relevante for å predikere 2025.
  \emph{Konsekvens:} Strukturelle brudd som betydelige prissjokk, atferdsendringer eller endringer i sortiment vil
  svekke prediksjonens gyldighet.
\item
  \textbf{Tilstrekkelig datakvalitet.} Datasettet antas å inneholde tilstrekkelig variasjon i nøkkeldimensjonene sesong,
  rabatt, region og produktkategori til å trene modellene. \emph{Konsekvens:} Simulerte data kan undervurdere virkelig
  støy og sjeldne hendelser, noe som kan gi mer optimistiske målinger enn hva som er realistisk i produksjon.
\item
  \textbf{Testsett som fremtidig periode.} 2025-dataene er holdt helt utenfor trening og brukes som testperiode.
  \emph{Konsekvens:} Evalueringen gir et realistisk bilde av prognosekvalitet under antakelsen om at 2025 ligner de
  foregående årene i mønster og struktur.
\end{enumerate}

\begin{center}\rule{0.5\linewidth}{0.5pt}\end{center}

\hypertarget{litteratur}{%
\subsection{2 Litteratur}\label{litteratur}}

I salgsprognose og etterspørselspredikering er to metodetradisjoner spesielt relevante: klassisk statistisk regresjon og
ensemble-basert maskinlæring. Begge er godt dokumenterte tilnærminger til prediktiv modellering i handels- og
forretningskontekster.

Multippel lineær regresjon er en veletablert statistisk metode for å predikere en kontinuerlig utfallsvariabel fra flere
forklaringsvariabler. IBM (u.å.-a) framhever at metoden er godt egnet for å avdekke mønstre i salgs- og innkjøpsdata og
hjelpe ledere med å forutsi etterspørselsperioder for produkter. GeeksforGeeks (2026a) understreker at lineær regresjon
er effektiv beregningskostnadsvis og gir et solid utgangspunkt for modellsammenlikning. Et fellestrekk i litteraturen er
at metoden krever at flere antagelser er oppfylt: residualene skal følge normalfordelingen, variansen skal være konstant
over alle prediktornivåer (homoskedastisitet), og det skal ikke være perfekt multikollinearitet mellom
forklaringsvariablene (IBM, u.å.-a; GeeksforGeeks, 2026a). Brudd på disse antagelsene svekker tolkningsvaliditeten, men
ikke nødvendigvis prediksjonskraften.

Random Forest representerer en nyere og mer fleksibel tilnærming. IBM (u.å.-b) og GeeksforGeeks (2026b) beskriver begge
at algoritmen er en ensemble-metode som kombinerer prediksjoner fra flere uavhengige beslutningstrær og dermed reduserer
variansen i prediksjoner sammenlignet med et enkelt tre. IBM (u.å.-b) framhever at Random Forest gir lett identifikasjon
av hvilke variabler som er viktigst for prediksjonen gjennom feature importance-rangeringer, og at metoden er robust mot
overfitting og håndterer manglende data godt. GeeksforGeeks (2026b) presiserer at algoritmen ikke krever normalisering
og er fleksibel for både klassifikasjon og regresjon.

Litteraturen gir ikke ett entydig svar på hvilken modell som er best for salgsprognose --- det avhenger av datastruktur,
evalueringsmål og krav til tolkbarhet. Begge metodene er vanlige valg, og valget av evalueringsmetrikk påvirker
konklusjonen. RMSE og MAPE er anerkjente metrikker for å sammenligne prognosemodeller, men de vektlegger henholdsvis
absolutt og relativ feil ulikt (GeeksforGeeks, 2026a). At de to metrikkene ikke alltid peker på samme vinner, er en
kjent utfordring i praktisk prognosearbeid.

De tilgjengelige kildene for dette prosjektet er webbaserte oppslagsverk fra IBM Think Topics og GeeksforGeeks, ikke
fagfellevurderte vitenskapelige artikler. Kildene gir likevel faglig solide introduksjoner til de valgte metodene og
tilstrekkelig grunnlag for metodevalgene i dette prosjektet. Supplering med fagfellevurderte studier --- for eksempel om
maskinlæring for etterspørselspredikering i dagligvarehandel --- bør vurderes dersom kravene til akademisk dybde øker.

\begin{center}\rule{0.5\linewidth}{0.5pt}\end{center}

\hypertarget{teori}{%
\subsection{3 Teori}\label{teori}}

Dette kapitlet gir det faglige grunnlaget for metodevalg, analyse og tolkning senere i rapporten. Notasjon: når en
kolonne eller feature fra modellmatrisen refereres direkte, skrives navnet i inline kode, for eksempel
\splitcode{Discount}, \splitcode{Region\_North} eller \splitcode{dayofmonth}. I Sammendrag og Abstract brukes vanlig
prosa («rabatt», «discount», «kalendervariabler») for å holde oppsummeringene tilgjengelige for et bredere publikum.
\emph{Forklaringsvariabel} brukes som samlebegrep for alle input i modellene, mens \emph{prediktor}, \emph{signal} og
\emph{feature} opptrer som kontekstuelle synonymer der variasjonen understøtter lesbarheten.

\hypertarget{multippel-lineuxe6r-regresjon}{%
\subsubsection{3.1 Multippel lineær regresjon}\label{multippel-lineuxe6r-regresjon}}

Multippel lineær regresjon er en supervisert læringsmetode som modellerer forholdet mellom en avhengig variabel og to
eller flere uavhengige forklaringsvariabler. Modellen uttrykkes som:

\(\hat{y} = \beta_0 + \beta_1 x_1 + \beta_2 x_2 + \cdots + \beta_n x_n\)

der \(\hat{y}\) er predikert verdi, \(\beta_0\) er konstantleddet og \(\beta_1, \ldots, \beta_n\) er
regresjonskoeffisientene som uttrykker endringen i \(\hat{y}\) per enhet endring i den tilhørende forklaringsvariabelen,
alt annet likt (GeeksforGeeks, 2026a).

Koeffisientene estimeres med \textbf{Ordinary Least Squares (OLS)}, som minimerer summen av kvadrerte residualer:

\(\sum_{i=1}^{n}(y_i - \hat{y}_i)^2\)

For at estimatene skal være gyldige krever metoden at fem antagelser er oppfylt (IBM, u.å.-a; GeeksforGeeks, 2026a): (1)
linearitet mellom avhengig og uavhengige variabler, (2) uavhengige residualer, (3) homoskedastisitet --- konstant
varians i residualene på tvers av alle prediktornivåer, (4) normalfordelte residualer, og (5) ingen perfekt
\textbf{multikollinearitet} mellom forklaringsvariablene. Multikollinearitet er særlig relevant i dette prosjektet fordi
one-hot-kodede dummyvariabler kan skape høy innbyrdes korrelasjon og gjøre tolkning av enkeltkoeffisienter usikker.

Lineær regresjon brukes her som benchmarkmodell. Dens styrke er tolkbarhet --- koeffisientene gir direkte innsikt i
retning og størrelse av variabeleffekter. Dens begrensning er at den forutsetter et lineært mønster i dataene.

\hypertarget{random-forest-regressor}{%
\subsubsection{3.2 Random Forest Regressor}\label{random-forest-regressor}}

Random Forest er en ensemble-metode basert på et sett med beslutningstrær som samlet gir mer stabile prediksjoner enn
ett enkelt tre (IBM, u.å.-b). Algoritmen hviler på to grunnprinsipper:

\textbf{Bootstrap aggregation (bagging):} Hvert tre trenes på et tilfeldig utvalg med tilbakelegging fra
treningsdataene. Dermed ser hvert tre en litt annerledes versjon av datasettet, og trærne korrelerer mindre med
hverandre (GeeksforGeeks, 2026b).

\textbf{Feature randomness:} Ved hver splittingsbeslutning i et tre velges kun et tilfeldig subsett av features til
vurdering --- ikke alle tilgjengelige variabler. Dette reduserer korrelasjon mellom trærne ytterligere og gjør ensemblet
mer robust mot overfitting (IBM, u.å.-b).

For regresjonsoppgaver er den endelige prediksjonen gjennomsnittet av alle trærnes individuelle prediksjoner. Dette
reduserer overordnet varians og demper effekten av støy i enkelttrær.

\textbf{Feature importance} beregnes som gjennomsnittlig reduksjon i MSE (Mean Decrease in Impurity) på tvers av alle
trær og splittinger for en gitt variabel. En variabel som konsekvent reduserer prediksjonsfeilen i mange trær, rangeres
høyt (IBM, u.å.-b).

Modellens oppførsel styres av hyperparametere som \splitcode{n\_estimators} (antall trær), \splitcode{max\_depth}
(maksimal dybde per tre), \splitcode{min\_samples\_leaf} (minste antall observasjoner i et løvnode) og
\splitcode{max\_features} (antall features vurdert per split). Disse tilpasses gjennom hyperparametertuning mot et
valideringssett.

\hypertarget{evalueringsmetrikker}{%
\subsubsection{3.3 Evalueringsmetrikker}\label{evalueringsmetrikker}}

To metrikker brukes for å evaluere prognosemodellene i dette prosjektet.

\textbf{RMSE (Root Mean Squared Error)} måler den gjennomsnittlige størrelsen på prediksjonsfeilen i samme enhet som
utfallsvariabelen:

\(\text{RMSE} = \sqrt{\frac{1}{n}\sum_{i=1}^{n}(y_i - \hat{y}_i)^2}\)

Kvadreringen medfører at store avvik vektes tyngre enn små. RMSE brukes som primær metrikk i dette prosjektet fordi
absolutt presisjon --- å treffe riktig salgsvolum --- er det mest relevante kravet for innkjøp og lagerstyring.

\textbf{MAPE (Mean Absolute Percentage Error)} måler det gjennomsnittlige prosentvise avviket:

\(\text{MAPE} = \frac{1}{n}\sum_{i=1}^{n}\left|\frac{y_i - \hat{y}_i}{y_i}\right| \times 100\)

MAPE er skalanøytral og enklere å kommunisere, men er ustabil når faktiske verdier er nær null (divisjon mot null).
Metrikken brukes som sekundær metrikk.

Når RMSE og MAPE peker på ulike vinnere, skyldes det at en modell kan ha lavt absolutt avvik på store volumer (lavt
RMSE) uten å treffe proporsjonalt godt på små volumer (høy MAPE). Begge metrikkene er derfor nødvendige for å forstå
modelloppførselen på tvers av ulike salgsnivåer (GeeksforGeeks, 2026a).

\hypertarget{feature-engineering-og-dataoppsett}{%
\subsubsection{3.4 Feature engineering og dataoppsett}\label{feature-engineering-og-dataoppsett}}

\textbf{Feature engineering} er prosessen med å utlede nye variabler fra rådata for å gjøre mønstre tilgjengelige for
modellene. I dette prosjektet er kalendervariablene \splitcode{year}, \splitcode{month}, \splitcode{quarter},
\splitcode{weekofyear}, \splitcode{dayofweek}, \splitcode{dayofmonth} og \splitcode{is\_weekend} avledet fra den
opprinnelige datovariabelen. Disse variablene fanger opp sesong-, uke- og kvartalsmønstre som ikke direkte kan leses fra
en rådate (GeeksforGeeks, 2026a).

\textbf{One-hot encoding} konverterer kategoriske variabler --- som Region, Category og Sub-Category --- til binære
dummyvariabler. For lineær regresjon er dette nødvendig fordi modellen krever numerisk input. Random Forest bruker her
samme kodede matrise for konsistens i sammenligning (IBM, u.å.-a).

\textbf{Data leakage} oppstår når variabler som ikke ville vært tilgjengelige på prediksjonstidspunktet inkluderes i
modellen. I dette prosjektet er \splitcode{Profit} ekskludert fordi den kun er kjent etter at salget er gjennomført.

\textbf{Tidsbasert oppsplitting} er valgt fremfor tilfeldig oppsplitting. Treningsdata er 2022--2024 og testdata er
2025. Tilfeldig oppsplitting ville tillate fremtidige observasjoner å inngå i treningen, noe som gir kunstig god ytelse
og ikke reflekterer reell prediksjon fremover i tid. I hyperparametertuningen av Random Forest brukes 2024 i tillegg som
intern valideringsperiode, mens lineær regresjon beholdes med det fulle treningsgrunnlaget.

\begin{center}\rule{0.5\linewidth}{0.5pt}\end{center}

\hypertarget{casebeskrivelse}{%
\subsection{4 Casebeskrivelse}\label{casebeskrivelse}}

\hypertarget{dagligvare-og-beslutningssituasjonen}{%
\subsubsection{4.1 Dagligvare og beslutningssituasjonen}\label{dagligvare-og-beslutningssituasjonen}}

Dagligvare er en simulert dagligvarekjede som opererer på tvers av fem regioner: West, East, Central, South og North.
Kjeden fører syv produktkategorier -- Snacks, Eggs/Meat/Fish, Fruits \& Veggies, Bakery, Beverages, Food Grains og Oil
\& Masala -- fordelt på 23 subkategorier. Datasettet som danner analysegrunnlaget inneholder 9 994 daglige
salgstransaksjoner fra perioden 2022--2025, fordelt på 24 byer.

Dagligvares operative planlegging er avhengig av pålitelige etterspørselsprognoser. Salget varierer med sesong, rabatter
og regionale forhold, og konsekvensene av dårlige prognoser er direkte synlige i driften: for høye bestillinger gir
svinn og bundet kapital, mens for lave bestillinger fører til utsolgte varer og tapte inntekter. Behovet for mer presise
prognoser er dermed sentralt for innkjøp, lager, kampanjevurdering og ressursplanlegging.

Figur 4.1 viser gjennomsnittlig salgsnivå per produktkategori. Eggs, Meat \& Fish skiller seg ut med det høyeste
gjennomsnittlige salgsnivået i datasettet.

\begin{figure}
\centering
\includegraphics[width=0.8\textwidth,height=\textheight]{../006 analysis/aktiviteter/07_eksplorativ_analyse_og_visualisering/fig_sales_per_category.png}
\caption{Figur 4.1 Gjennomsnittlig salgsnivå per produktkategori.}
\end{figure}

\hypertarget{historisk-salgsutvikling}{%
\subsubsection{4.2 Historisk salgsutvikling}\label{historisk-salgsutvikling}}

Gjennomsnittlig daglig salg for hele perioden 2022--2025 er 1 497 med et spenn fra 500 til 2 500 og standardavvik på
578. Salgsnivået er stabilt over tid: gjennomsnittsnivået i treningsperioden 2022--2024 er 1 493, mot 1 503 i
testperioden 2025. Ingen tydelig veksttrend er synlig i perioden -- nivået holder seg relativt konstant med lokal
variabilitet.

Figur 4.2 viser salgsforløpet over tid med trenings- og testperioden markert.

\begin{figure}
\centering
\includegraphics[width=0.8\textwidth,height=\textheight]{../006 analysis/aktiviteter/07_eksplorativ_analyse_og_visualisering/fig_sales_over_tid_train_test.png}
\caption{Figur 4.2 Daglig salg over perioden 2022--2025, med treningsperiode (2022--2024) og testperiode (2025) vist.}
\end{figure}

\hypertarget{sesongmuxf8nster-i-salget}{%
\subsubsection{4.3 Sesongmønster i salget}\label{sesongmuxf8nster-i-salget}}

Månedsmønsteret viser at oktober er måneden med det høyeste gjennomsnittlige salgsnivået, mens juni er den svakeste
måneden. Mønsteret er konsistent mellom trenings- og testperioden, noe som tyder på et stabilt sesongmønster i
datasettet. Fjerde kvartal fremstår samlet som den sterkeste perioden, mens andre og tredje kvartal er relativt svakere.

Figur 4.3 viser gjennomsnittlig salg per måned for henholdsvis trenings- og testperioden.

\begin{figure}
\centering
\includegraphics[width=0.8\textwidth,height=\textheight]{../006 analysis/aktiviteter/07_eksplorativ_analyse_og_visualisering/fig_sales_per_month_split.png}
\caption{Figur 4.3 Gjennomsnittlig daglig salg per måned for treningsperioden (2022--2024) og testperioden (2025).}
\end{figure}

\hypertarget{utfordringer-duxe5rlige-prognoser-medfuxf8rer-i-bedriften}{%
\subsubsection{4.4 Utfordringer dårlige prognoser medfører i
bedriften}\label{utfordringer-duxe5rlige-prognoser-medfuxf8rer-i-bedriften}}

For Dagligvare forsterkes planleggingsrisikoen av et sortiment som spenner fra ferske varer med kort holdbarhet til
tørrvarer med lengre holdbarhetsperiode. Overskuddsbestillinger fører til direkte svinn og bundet kapital --
konsekvenser som er særlig kostbare i kategorier som Eggs, Meat \& Fish og Fruits \& Veggies. Underestimering gir tomme
hyller, tapte inntekter og risiko for å svekke kundelojaliteten.

Sesongvariasjonen i Figur 4.3 viser at salget svinger markant gjennom året -- med topp i oktober og bunnpunkt i juni.
Kombinert med rabattavhengighet gjør dette at enkle tommelfingerregler for innkjøp gir dårlige resultater i perioder med
høy kampanjeaktivitet eller sesongtopper.

Nøyaktige etterspørselsprognoser for 2025 kan gi Dagligvare et bedre grunnlag for innkjøpsplanlegging og lagerstyring,
støtte vurderingen av kampanjeeffekter og bidra til mer presis ressursplanlegging i perioder med høy og lav
etterspørsel.

\begin{center}\rule{0.5\linewidth}{0.5pt}\end{center}

\hypertarget{metode-og-data}{%
\subsection{5 Metode og data}\label{metode-og-data}}

\hypertarget{metode}{%
\subsubsection{5.1 Metode}\label{metode}}

Prosjektet er en kvantitativ, prediktiv studie der historiske salgsdata fra én simulert dagligvarekjede brukes til å
bygge og evaluere prognosemodeller. Datagrunnlaget er ikke samlet inn, men stilt til disposisjon som en del av
prosjektets faglige rammer. Problemstillingen besvares gjennom en trinnvis analytisk prosess: datasettet ble renset og
gjort klart for modellering, deretter ble relevante variabler valgt og nye kalendervariabler utledet, og til slutt ble
datasettet delt i en treningsperiode og en testperiode basert på tid.

Lineær regresjon ble valgt som benchmark-modell fordi metoden er tolkbar og gir et stabilt sammenligningsgrunnlag.
Random Forest Regressor ble valgt som alternativmodell fordi den kan fange opp ikke-lineære mønstre og gir en naturlig
rangering av variablenes prediksjonsverdi. Random Forest ble i tillegg tunet ved å trene på 2022--2023 og validere på
2024, slik at hyperparametere ble valgt uten å bruke testdataene. Alle tre modellspor -- lineær regresjon, Random Forest
baseline og tuned Random Forest -- ble evaluert på 2025-data.

Evalueringen benytter RMSE som primær metrikk fordi absolutt presisjon er mest relevant for innkjøp og lagerstyring, og
MAPE som sekundær metrikk for å gi et relativt bilde av prognosefeilen (se kap. 3.3). Modellene sammenlignes samlet for
hele 2025 og per måned, og resultatene tolkes videre etter kvartal, rabattnivå, region og salgsnivå for å gi praktisk
beslutningsstøtte til Dagligvare.

\hypertarget{data}{%
\subsubsection{5.2 Data}\label{data}}

Datagrunnlaget er filen \splitcode{Dagligvare\_Dataset.csv}, som inneholder 9 994 daglige salgstransaksjoner fra en
simulert dagligvarekjede. Rådata dekker opprinnelig perioden 2015--2018 og er remappet med en syvårig
kalenderforskyvning til prosjektperioden 2022--2025. Datasettet har 11 råkolonner og ingen manglende verdier eller
dubletter. To ulike datoformater i kildedata -- 4 042 rader med \splitcode{dd-mm-yyyy} og 5 952 rader med
\splitcode{mm/dd/yyyy} -- ble standardisert til ISO-format under rensingen.

Figur 5.1 viser fordelingen av datatyper i rådatasettet og underbygger valget av forbehandling og feature engineering.

\begin{figure}
\centering
\includegraphics[width=0.8\textwidth,height=\textheight]{../006 analysis/aktiviteter/01_dataforstaelse_og_variabler/fig_datatype_fordeling.png}
\caption{Figur 5.1 Fordeling av datatyper i rådatasettet. Fordelingen viser hvor mange kolonner som er numeriske,
kategoriske og datobaserte, og underbygger valg av forbehandling og feature engineering.}
\end{figure}

Målvariabelen er \splitcode{Sales} (heltall, spenn 500--2 500). Forklaringsvariablene som inngår i modellene er
\splitcode{Discount} (desimaltall, 0,10--0,35) og de kategoriske variablene \splitcode{Category},
\splitcode{Sub Category}, \splitcode{City} og \splitcode{Region}. Fra \splitcode{Order Date} er det i tillegg utledet
sju kalendervariabler: \splitcode{year}, \splitcode{month}, \splitcode{quarter}, \splitcode{weekofyear},
\splitcode{dayofweek}, \splitcode{dayofmonth} og \splitcode{is\_weekend}. Fire variabler ble ekskludert:
\splitcode{Profit} fordi den er en lekkasjevariabel som kun er kjent etter gjennomført salg, \splitcode{State} fordi
kolonnen er konstant i datasettet, og \splitcode{Order ID} og \splitcode{Customer Name} fordi de ikke har generaliserbar
prediksjonsverdi. Kategoriske variabler ble one-hot-encoded til binære dummyvariabler, slik at den endelige
modellmatrisen inneholder 67 features.

Tabell 5.1 oppsummerer de 11 kolonnene i rådatasettet med datatype, manglende andel og anbefaling for videre bruk.

\begin{longtable}[]{@{}
  >{\raggedright\arraybackslash}p{(\columnwidth - 10\tabcolsep) * \real{0.1667}}
  >{\raggedright\arraybackslash}p{(\columnwidth - 10\tabcolsep) * \real{0.1667}}
  >{\raggedright\arraybackslash}p{(\columnwidth - 10\tabcolsep) * \real{0.1667}}
  >{\raggedright\arraybackslash}p{(\columnwidth - 10\tabcolsep) * \real{0.1667}}
  >{\raggedright\arraybackslash}p{(\columnwidth - 10\tabcolsep) * \real{0.1667}}
  >{\raggedright\arraybackslash}p{(\columnwidth - 10\tabcolsep) * \real{0.1667}}@{}}
\caption{Tabell 5.1 Variabeloversikt med datatype, manglende andel og anbefaling for videre bruk.}\tabularnewline
\toprule\noalign{}
\begin{minipage}[b]{\linewidth}\raggedright
Variabel
\end{minipage} & \begin{minipage}[b]{\linewidth}\raggedright
Datatype
\end{minipage} & \begin{minipage}[b]{\linewidth}\raggedright
Manglende \%
\end{minipage} & \begin{minipage}[b]{\linewidth}\raggedright
Unike
\end{minipage} & \begin{minipage}[b]{\linewidth}\raggedright
Anbefaling
\end{minipage} & \begin{minipage}[b]{\linewidth}\raggedright
\splitword{Begrunnelse}
\end{minipage} \\
\midrule\noalign{}
\endfirsthead
\toprule\noalign{}
\begin{minipage}[b]{\linewidth}\raggedright
Variabel
\end{minipage} & \begin{minipage}[b]{\linewidth}\raggedright
Datatype
\end{minipage} & \begin{minipage}[b]{\linewidth}\raggedright
Manglende \%
\end{minipage} & \begin{minipage}[b]{\linewidth}\raggedright
Unike
\end{minipage} & \begin{minipage}[b]{\linewidth}\raggedright
Anbefaling
\end{minipage} & \begin{minipage}[b]{\linewidth}\raggedright
\splitword{Begrunnelse}
\end{minipage} \\
\midrule\noalign{}
\endhead
\bottomrule\noalign{}
\endlastfoot
Sales & int64 & 0,0 & 1 989 & target & Målvariabel for prognose. \\
Category & str & 0,0 & 7 & inkluder & Kategorisk variabel med håndterbar kardinalitet. \\
City & str & 0,0 & 24 & inkluder & Kategorisk variabel med håndterbar kardinalitet. \\
Discount & float64 & 0,0 & 26 & inkluder & Numerisk variabel egner seg direkte for modellering. \\
Region & str & 0,0 & 5 & inkluder & Kategorisk variabel med håndterbar kardinalitet. \\
State & str & 0,0 & 1 & ekskluder & Kolonnen er konstant i datasettet og gir ingen \splitword{forklaringskraft.} \\
Sub Category & str & 0,0 & 23 & inkluder & Kategorisk variabel med håndterbar kardinalitet. \\
Order Date & str & 0,0 & 1 236 & vurder & Middels/høy kardinalitet, vurder koding og nytte. \\
Profit & float64 & 0,0 & 8 380 & vurder & Kan være informativ, men bør sjekkes for lekkasje før modellering. \\
Customer Name & str & 0,0 & 50 & ekskluder & Navn-kolonne har høy kardinalitet og høy risiko for
\splitword{overtilpasning.} \\
Order ID & str & 0,0 & 9 994 & ekskluder & ID-variabel gir vanligvis lite generaliserbar
\splitword{prediksjonsverdi.} \\
\end{longtable}

Tabell 5.2 dokumenterer feature engineering-oppsettet: hvilke input-kolonner som beholdes, transformeres eller
ekskluderes, og hvilke output-kolonner de gir.

\begin{longtable}[]{@{}
  >{\raggedright\arraybackslash}p{(\columnwidth - 6\tabcolsep) * \real{0.2500}}
  >{\raggedright\arraybackslash}p{(\columnwidth - 6\tabcolsep) * \real{0.2500}}
  >{\raggedright\arraybackslash}p{(\columnwidth - 6\tabcolsep) * \real{0.2500}}
  >{\raggedright\arraybackslash}p{(\columnwidth - 6\tabcolsep) * \real{0.2500}}@{}}
\caption{Tabell 5.2 Feature engineering-oppsett: input-kolonne, handling og resulterende output-kolonne med
begrunnelse.}\tabularnewline
\toprule\noalign{}
\begin{minipage}[b]{\linewidth}\raggedright
Input-kolonne
\end{minipage} & \begin{minipage}[b]{\linewidth}\raggedright
Handling
\end{minipage} & \begin{minipage}[b]{\linewidth}\raggedright
Output-kolonne
\end{minipage} & \begin{minipage}[b]{\linewidth}\raggedright
\splitword{Begrunnelse}
\end{minipage} \\
\midrule\noalign{}
\endfirsthead
\toprule\noalign{}
\begin{minipage}[b]{\linewidth}\raggedright
Input-kolonne
\end{minipage} & \begin{minipage}[b]{\linewidth}\raggedright
Handling
\end{minipage} & \begin{minipage}[b]{\linewidth}\raggedright
Output-kolonne
\end{minipage} & \begin{minipage}[b]{\linewidth}\raggedright
\splitword{Begrunnelse}
\end{minipage} \\
\midrule\noalign{}
\endhead
\bottomrule\noalign{}
\endlastfoot
Sales & behold & Sales & Målvariabel for prognose. \\
Discount & behold & Discount & Numerisk \splitword{forklaringsvariabel.} \\
Category & behold & Category & Kategorisk feature med håndterbar kardinalitet. \\
Sub Category & behold & Sub Category & Kategorisk feature med håndterbar kardinalitet. \\
City & behold & City & Kategorisk feature med håndterbar kardinalitet. \\
Region & behold & Region & Kategorisk feature med håndterbar kardinalitet. \\
Order Date & behold & Order Date & Beholdes for sporbarhet og tidsbasert splitt. \\
Order Date & transformer & year & Utledet kalenderår etter remapping. \\
Order Date & transformer & month & Utledet måned. \\
Order Date & transformer & quarter & Utledet kvartal. \\
Order Date & transformer & weekofyear & Utledet uke i året. \\
Order Date & transformer & dayofweek & Utledet ukedag. \\
Order Date & transformer & dayofmonth & Utledet dag i måneden. \\
Order Date & transformer & is\_weekend & Indikator for helg. \\
Order ID & ekskluder & -- & ID-kolonne uten generaliserbar \splitword{prediksjonsverdi.} \\
Customer Name & ekskluder & -- & Navn-kolonne med høy risiko for \splitword{overtilpasning.} \\
Profit & ekskluder & -- & Ekskludert som potensiell \splitword{lekkasjevariabel.} \\
State & ekskluder & -- & Kolonnen er konstant i datasettet og gir ingen \splitword{forklaringskraft.} \\
\end{longtable}

Tabell 5.3 dokumenterer datarensingen fra 9 994 rader inn til 9 994 rader ut, uten at rader ble fjernet eller manglende
verdier oppdaget.

\begin{longtable}[]{@{}
  >{\raggedright\arraybackslash}p{(\columnwidth - 4\tabcolsep) * \real{0.3333}}
  >{\raggedright\arraybackslash}p{(\columnwidth - 4\tabcolsep) * \real{0.3333}}
  >{\raggedright\arraybackslash}p{(\columnwidth - 4\tabcolsep) * \real{0.3333}}@{}}
\caption{Tabell 5.3 Datarensing -- målepunkter, verdier og kommentarer som dokumenterer kvalitetskontrollen før
modellering.}\tabularnewline
\toprule\noalign{}
\begin{minipage}[b]{\linewidth}\raggedright
Målepunkt
\end{minipage} & \begin{minipage}[b]{\linewidth}\raggedright
Verdi
\end{minipage} & \begin{minipage}[b]{\linewidth}\raggedright
Kommentar
\end{minipage} \\
\midrule\noalign{}
\endfirsthead
\toprule\noalign{}
\begin{minipage}[b]{\linewidth}\raggedright
Målepunkt
\end{minipage} & \begin{minipage}[b]{\linewidth}\raggedright
Verdi
\end{minipage} & \begin{minipage}[b]{\linewidth}\raggedright
Kommentar
\end{minipage} \\
\midrule\noalign{}
\endhead
\bottomrule\noalign{}
\endlastfoot
\splitword{antall\_rader\_inn} & 9 994 & Observasjoner i rådata \\
\splitword{antall\_kolonner\_inn} & 11 & Variabler i rådata \\
\splitword{manglende\_verdier\_inn} & 0 & Totalt manglende i rådata \\
dubletter\_inn & 0 & Eksakte dublettrader i rådata \\
\splitword{datoformat\_dd\_mm\_yyyy} & 4 042 & Tolket fra verdier med bindestrek \\
\splitword{datoformat\_mm\_dd\_yyyy} & 5 952 & Tolket fra verdier med skråstrek \\
\splitword{datoformat\_annet} & 0 & Verdier uten kjent mønster \\
\splitword{ugyldige\_datoer\_etter\_tolkning} & 0 & Datoer som ikke lot seg parse \\
\splitword{opprinnelig\_periode\_start} & 2015-01-02 & Tidligste dato i original dataserie \\
\splitword{opprinnelig\_periode\_slutt} & 2018-12-30 & Siste dato i original dataserie \\
år\_forskyvning & 7 & \splitword{Kalenderforskyvning} brukt for prosjektets \splitword{arbeidsgrunnlag} \\
\splitword{prosjekt\_periode\_start} & 2022-01-02 & Tidligste dato etter remapping \\
\splitword{prosjekt\_periode\_slutt} & 2025-12-30 & Siste dato etter remapping \\
\splitword{rader\_fjernet\_manglende\_kritisk\_felt} & 0 & Rader fjernet etter kontroll av dato og numeriske
nøkkelfelt \\
\splitword{dubletter\_fjernet} & 0 & Fjernet etter \splitword{standardisering} \\
\splitword{antall\_rader\_ut} & 9 994 & Observasjoner i renset datasett \\
\splitword{manglende\_verdier\_ut} & 0 & Totalt manglende etter rens \\
dubletter\_ut & 0 & Eksakte dublettrader etter rens \\
\end{longtable}

Fordelingen av kategoriske variabler viser at West-regionen utgjør 32 \% av observasjonene, East 28 \%, Central 23 \% og
South 16 \%. North-regionen er representert med bare én observasjon og benyttes ikke som separat analysesegment.
Produktkategoriene er jevnt fordelt med mellom 14 og 15 \% av observasjonene hver, noe som betyr at ingen enkelt
kategori dominerer datasettet.

Datasettet er delt tidsmessig slik at treningsdata dekker 2022--2024 og testdata dekker 2025. Tilfeldig splitt ble ikke
brukt, fordi det ville tillate fremtidige observasjoner å inngå i treningen og dermed gi kunstig god ytelse. Det
gjennomsnittlige salgsnivået er stabilt mellom periodene (1 493 i treningsdata og 1 503 i testdata), noe som styrker
antakelsen om at historiske mønstre er overførbare til 2025. Figur 5.2 viser fordelingen av daglige salgsverdier i
trenings- og testperioden.

\begin{figure}
\centering
\includegraphics[width=0.8\textwidth,height=\textheight]{../006 analysis/aktiviteter/07_eksplorativ_analyse_og_visualisering/fig_sales_fordeling_train_test.png}
\caption{Figur 5.2 Fordeling av daglige salgsverdier i trenings- og testperioden. Overlappet i fordelingene tyder på at
målvariabelen er stabil mellom periodene.}
\end{figure}

Tabell 5.4 oppsummerer antallsfordelingen mellom trenings- og testperioden.

\begin{longtable}[]{@{}
  >{\raggedright\arraybackslash}p{(\columnwidth - 6\tabcolsep) * \real{0.2500}}
  >{\raggedright\arraybackslash}p{(\columnwidth - 6\tabcolsep) * \real{0.2500}}
  >{\raggedright\arraybackslash}p{(\columnwidth - 6\tabcolsep) * \real{0.2500}}
  >{\raggedright\arraybackslash}p{(\columnwidth - 6\tabcolsep) * \real{0.2500}}@{}}
\caption{Tabell 5.4 Fordeling av trenings- og testdata etter tidsmessig splitt.}\tabularnewline
\toprule\noalign{}
\begin{minipage}[b]{\linewidth}\raggedright
Delmengde
\end{minipage} & \begin{minipage}[b]{\linewidth}\raggedright
År
\end{minipage} & \begin{minipage}[b]{\linewidth}\raggedright
Antall rader
\end{minipage} & \begin{minipage}[b]{\linewidth}\raggedright
Andel
\end{minipage} \\
\midrule\noalign{}
\endfirsthead
\toprule\noalign{}
\begin{minipage}[b]{\linewidth}\raggedright
Delmengde
\end{minipage} & \begin{minipage}[b]{\linewidth}\raggedright
År
\end{minipage} & \begin{minipage}[b]{\linewidth}\raggedright
Antall rader
\end{minipage} & \begin{minipage}[b]{\linewidth}\raggedright
Andel
\end{minipage} \\
\midrule\noalign{}
\endhead
\bottomrule\noalign{}
\endlastfoot
Treningsdata & 2022--2024 & 6 682 & \textasciitilde67 \% \\
Testdata & 2025 & 3 312 & \textasciitilde33 \% \\
\end{longtable}

\begin{center}\rule{0.5\linewidth}{0.5pt}\end{center}

\hypertarget{modellering}{%
\subsection{6 Modellering}\label{modellering}}

Prosjektet evaluerer tre modellspor for å predikere daglig salg i 2025. Benchmark-modellen er en multippel lineær
regresjon som gir et tolkbart sammenligningsgrunnlag. Random Forest inngår i to varianter: en baseline med
standardparametere som viser hva en utuned ensemble-modell leverer på det samme datagrunnlaget, og en tuned variant der
hyperparametere er valgt gjennom et strukturert søk. Sammenstillingen gir både et grunnlag for å måle gevinsten av
tuning og et tolkbart referansepunkt som er robust mot overfitting.

Tabell 6.1 oppsummerer de tre modellsporene med antall features, treningsrader, sentrale hyperparametre og rolle i
prosjektet.

\begin{longtable}[]{@{}
  >{\raggedright\arraybackslash}p{(\columnwidth - 10\tabcolsep) * \real{0.1667}}
  >{\raggedright\arraybackslash}p{(\columnwidth - 10\tabcolsep) * \real{0.1667}}
  >{\raggedright\arraybackslash}p{(\columnwidth - 10\tabcolsep) * \real{0.1667}}
  >{\raggedright\arraybackslash}p{(\columnwidth - 10\tabcolsep) * \real{0.1667}}
  >{\raggedright\arraybackslash}p{(\columnwidth - 10\tabcolsep) * \real{0.1667}}
  >{\raggedright\arraybackslash}p{(\columnwidth - 10\tabcolsep) * \real{0.1667}}@{}}
\caption{Tabell 6.1 Oversikt over de tre modellsporene: antall features, treningsrader, sentrale hyperparametre og rolle
i prosjektet.}\tabularnewline
\toprule\noalign{}
\begin{minipage}[b]{\linewidth}\raggedright
\splitword{Modellrolle}
\end{minipage} & \begin{minipage}[b]{\linewidth}\raggedright
Algoritme
\end{minipage} & \begin{minipage}[b]{\linewidth}\raggedright
\splitword{Treningsrader}
\end{minipage} & \begin{minipage}[b]{\linewidth}\raggedright
Features
\end{minipage} & \begin{minipage}[b]{\linewidth}\raggedright
Sentrale \splitword{hyperparametre}
\end{minipage} & \begin{minipage}[b]{\linewidth}\raggedright
Rolle i prosjektet
\end{minipage} \\
\midrule\noalign{}
\endfirsthead
\toprule\noalign{}
\begin{minipage}[b]{\linewidth}\raggedright
\splitword{Modellrolle}
\end{minipage} & \begin{minipage}[b]{\linewidth}\raggedright
Algoritme
\end{minipage} & \begin{minipage}[b]{\linewidth}\raggedright
\splitword{Treningsrader}
\end{minipage} & \begin{minipage}[b]{\linewidth}\raggedright
Features
\end{minipage} & \begin{minipage}[b]{\linewidth}\raggedright
Sentrale \splitword{hyperparametre}
\end{minipage} & \begin{minipage}[b]{\linewidth}\raggedright
Rolle i prosjektet
\end{minipage} \\
\midrule\noalign{}
\endhead
\bottomrule\noalign{}
\endlastfoot
benchmark lineær & \splitword{LinearRegression} & 6 682 & 67 & \splitcode{fit\_intercept=True}, ingen regularisering &
Tolkbar benchmark med eksplisitt \splitword{koeffisientretning.} \\
baseline RF & \splitword{RandomForestRegressor} & 6 682 & 67 & \splitcode{n\_estimators=200},
\splitcode{max\_depth=None}, \splitcode{min\_samples\_leaf=1}, \splitcode{max\_features=1.0} & RF-referanse uten tuning,
for \splitword{stabilitetssjekk} mot tunet modell. \\
tuned RF & \splitword{RandomForestRegressor} & 6 682 & 67 & \splitcode{n\_estimators=400}, \splitcode{max\_depth=10},
\splitcode{min\_samples\_leaf=4}, \splitcode{max\_features="sqrt"} & Anbefalt operativ modell etter tuning på
2024-validering. \\
\end{longtable}

Den multiple lineære regresjonen trenes som \splitcode{LinearRegression} med \splitcode{fit\_intercept=True} og uten
regularisering eller skalering. Modellen estimeres med OLS (se kap. 3.1) på 6 682 rader og 67 features -- det vil si
alle kalendervariabler, \splitcode{Discount} og de one-hot-kodede kategoriske variablene. Estimert konstantledd er
\(-3193{,}55\). Modellen er tolkbar gjennom sine koeffisienter, men forutsetter at salget kan beskrives som en lineær
kombinasjon av forklaringsvariablene, og den har ingen mekanisme for å fange opp interaksjoner eller ikke-lineære
effekter utover de som eksplisitt er kodet inn.

Random Forest Regressor trenes på samme treningsmatrise. Baselinen bruker standardkonfigurasjonen
\splitcode{n\_estimators=200}, \splitcode{max\_depth=None}, \splitcode{min\_samples\_leaf=1},
\splitcode{max\_features=1.0} (dvs. at alle 67 features vurderes ved hver split), \splitcode{bootstrap=True} og
\splitcode{random\_state=42}, jf. bagging og feature randomness i kap. 3.2. Den tunede varianten bygger på et
rutenettsøk der hver kandidat trenes på 2022--2023 (4 095 rader) og valideres på 2024 (2 587 rader). 2025-data inngår
ikke i modellutvelgelsen -- denne tidsmessige isolasjonen forhindrer datalekkasje mellom tuning og evaluering (jf. kap.
3.4 og 5.2). Vinnerkonfigurasjonen ble \splitcode{n\_estimators=400}, \splitcode{max\_depth=10},
\splitcode{min\_samples\_leaf=4} og \splitcode{max\_features='sqrt'} (som gir \(\lfloor\sqrt{67}\rfloor = 8\) features
per split), valgt med validerings-RMSE som primærkriterium. Vinnerkandidaten oppnådde validerings-RMSE \(577{,}27\) og
validerings-MAPE \(43{,}56\%\), mot baseline-kandidatens \(590{,}30\) og \(44{,}22\%\) på samme valideringssett. Den
tunede modellen retrenes deretter på hele treningsperioden (2022--2024) før evaluering på 2025. Alle tre modellene
genererer prognoser på testperioden 2025 (3 312 rader), og det er disse prognosene som danner grunnlaget for analyse og
resultat i de neste kapitlene.

Tabell 6.2 viser de fem beste parameterkombinasjonene fra rutenettsøket, sortert etter RMSE på 2024-valideringen og med
delta mot baseline-kandidaten.

\begin{longtable}[]{@{}
  >{\raggedright\arraybackslash}p{(\columnwidth - 14\tabcolsep) * \real{0.1250}}
  >{\raggedright\arraybackslash}p{(\columnwidth - 14\tabcolsep) * \real{0.1250}}
  >{\raggedright\arraybackslash}p{(\columnwidth - 14\tabcolsep) * \real{0.1250}}
  >{\raggedright\arraybackslash}p{(\columnwidth - 14\tabcolsep) * \real{0.1250}}
  >{\raggedright\arraybackslash}p{(\columnwidth - 14\tabcolsep) * \real{0.1250}}
  >{\raggedright\arraybackslash}p{(\columnwidth - 14\tabcolsep) * \real{0.1250}}
  >{\raggedright\arraybackslash}p{(\columnwidth - 14\tabcolsep) * \real{0.1250}}
  >{\raggedright\arraybackslash}p{(\columnwidth - 14\tabcolsep) * \real{0.1250}}@{}}
\caption{Tabell 6.2 Topp 5 parameterkombinasjoner fra tuning av Random Forest, sortert etter RMSE på
2024-valideringen.}\tabularnewline
\toprule\noalign{}
\begin{minipage}[b]{\linewidth}\raggedright
Kandidat
\end{minipage} & \begin{minipage}[b]{\linewidth}\raggedright
\splitword{n\_estimators}
\end{minipage} & \begin{minipage}[b]{\linewidth}\raggedright
max\_depth
\end{minipage} & \begin{minipage}[b]{\linewidth}\raggedright
\splitword{min\_samples\_leaf}
\end{minipage} & \begin{minipage}[b]{\linewidth}\raggedright
\splitword{max\_features}
\end{minipage} & \begin{minipage}[b]{\linewidth}\raggedright
RMSE validering
\end{minipage} & \begin{minipage}[b]{\linewidth}\raggedright
MAPE validering (\%)
\end{minipage} & \begin{minipage}[b]{\linewidth}\raggedright
Delta RMSE \splitword{vs. baseline}
\end{minipage} \\
\midrule\noalign{}
\endfirsthead
\toprule\noalign{}
\begin{minipage}[b]{\linewidth}\raggedright
Kandidat
\end{minipage} & \begin{minipage}[b]{\linewidth}\raggedright
\splitword{n\_estimators}
\end{minipage} & \begin{minipage}[b]{\linewidth}\raggedright
max\_depth
\end{minipage} & \begin{minipage}[b]{\linewidth}\raggedright
\splitword{min\_samples\_leaf}
\end{minipage} & \begin{minipage}[b]{\linewidth}\raggedright
\splitword{max\_features}
\end{minipage} & \begin{minipage}[b]{\linewidth}\raggedright
RMSE validering
\end{minipage} & \begin{minipage}[b]{\linewidth}\raggedright
MAPE validering (\%)
\end{minipage} & \begin{minipage}[b]{\linewidth}\raggedright
Delta RMSE \splitword{vs. baseline}
\end{minipage} \\
\midrule\noalign{}
\endhead
\bottomrule\noalign{}
\endlastfoot
rf\_tune\_30 (vinner) & 400 & 10 & 4 & sqrt & 577,27 & 43,56 & −13,04 \\
rf\_tune\_28 & 400 & 10 & 2 & sqrt & 577,29 & 43,57 & −13,02 \\
rf\_tune\_10 & 200 & 10 & 2 & sqrt & 577,36 & 43,58 & −12,94 \\
rf\_tune\_26 & 400 & 10 & 1 & sqrt & 577,65 & 43,60 & −12,65 \\
rf\_tune\_12 & 200 & 10 & 4 & sqrt & 577,81 & 43,61 & −12,50 \\
\end{longtable}

\begin{center}\rule{0.5\linewidth}{0.5pt}\end{center}

\hypertarget{analyse}{%
\subsection{7 Analyse}\label{analyse}}

Samlet for 2025 leverer den tunede Random Forest-modellen både lavest RMSE og lavest MAPE (jf. Tabell 8.1), marginalt
foran benchmark lineær og tydeligere foran baseline Random Forest. Gapet mellom modellene er lite i absolutt forstand --
tuned Random Forest forbedrer RMSE med omtrent \(11\) enheter sammenlignet med baseline Random Forest og \(2\) enheter
sammenlignet med benchmark lineær. Tuningeffekten tolkes som en kombinasjon av dybdebegrensningen
(\splitcode{max\_depth=10}) og redusert \splitcode{max\_features}, som demper variansen i de enkelte trærne og gir mer
stabile ensemble-prediksjoner på nye data (jf. kap. 3.2).

Månedsnivået nyanserer bildet vesentlig (jf. Tabell 7.1). Tuned Random Forest vinner RMSE i 11 av 12 måneder, men vinner
MAPE i bare 3 måneder. Baseline Random Forest vinner ikke en eneste måned på RMSE, men vinner MAPE i 6 av 12 måneder.
RMSE og MAPE peker på samme vinner bare i juli og september. Mønsteret kan forklares med at absolutt feil og prosentfeil
vekter observasjoner ulikt: RMSE straffer store avvik hardt og favoriserer dermed modellen som er mest treffsikker på de
største salgsmålene, mens MAPE er mer sensitiv i måneder med mange lavvolumrader der små avvik gir utslag i høye
prosenttall.

Tabell 7.1 viser eksakte månedlige RMSE- og MAPE-verdier per modell.

\begin{longtable}[]{@{}
  >{\raggedright\arraybackslash}p{(\columnwidth - 12\tabcolsep) * \real{0.1429}}
  >{\raggedright\arraybackslash}p{(\columnwidth - 12\tabcolsep) * \real{0.1429}}
  >{\raggedright\arraybackslash}p{(\columnwidth - 12\tabcolsep) * \real{0.1429}}
  >{\raggedright\arraybackslash}p{(\columnwidth - 12\tabcolsep) * \real{0.1429}}
  >{\raggedright\arraybackslash}p{(\columnwidth - 12\tabcolsep) * \real{0.1429}}
  >{\raggedright\arraybackslash}p{(\columnwidth - 12\tabcolsep) * \real{0.1429}}
  >{\raggedright\arraybackslash}p{(\columnwidth - 12\tabcolsep) * \real{0.1429}}@{}}
\caption{Tabell 7.1 Månedlig RMSE og MAPE per modell i 2025.}\tabularnewline
\toprule\noalign{}
\begin{minipage}[b]{\linewidth}\raggedright
Måned
\end{minipage} & \begin{minipage}[b]{\linewidth}\raggedright
RMSE lineær
\end{minipage} & \begin{minipage}[b]{\linewidth}\raggedright
RMSE baseline RF
\end{minipage} & \begin{minipage}[b]{\linewidth}\raggedright
RMSE tuned RF
\end{minipage} & \begin{minipage}[b]{\linewidth}\raggedright
MAPE lineær (\%)
\end{minipage} & \begin{minipage}[b]{\linewidth}\raggedright
MAPE baseline RF (\%)
\end{minipage} & \begin{minipage}[b]{\linewidth}\raggedright
MAPE tuned RF (\%)
\end{minipage} \\
\midrule\noalign{}
\endfirsthead
\toprule\noalign{}
\begin{minipage}[b]{\linewidth}\raggedright
Måned
\end{minipage} & \begin{minipage}[b]{\linewidth}\raggedright
RMSE lineær
\end{minipage} & \begin{minipage}[b]{\linewidth}\raggedright
RMSE baseline RF
\end{minipage} & \begin{minipage}[b]{\linewidth}\raggedright
RMSE tuned RF
\end{minipage} & \begin{minipage}[b]{\linewidth}\raggedright
MAPE lineær (\%)
\end{minipage} & \begin{minipage}[b]{\linewidth}\raggedright
MAPE baseline RF (\%)
\end{minipage} & \begin{minipage}[b]{\linewidth}\raggedright
MAPE tuned RF (\%)
\end{minipage} \\
\midrule\noalign{}
\endhead
\bottomrule\noalign{}
\endlastfoot
2025-01 & 586,09 & 597,82 & 585,73 & 46,62 & 46,51 & 46,52 \\
2025-02 & 577,96 & 584,01 & 576,74 & 42,50 & 42,19 & 42,72 \\
2025-03 & 552,52 & 565,11 & 550,33 & 42,14 & 42,86 & 42,30 \\
2025-04 & 581,28 & 576,18 & 575,63 & 44,76 & 43,04 & 43,85 \\
2025-05 & 591,70 & 594,57 & 589,74 & 45,46 & 44,36 & 45,00 \\
2025-06 & 560,16 & 556,89 & 552,25 & 43,81 & 41,87 & 42,79 \\
2025-07 & 569,61 & 582,87 & 569,23 & 42,90 & 42,67 & 42,33 \\
2025-08 & 586,60 & 601,17 & 586,65 & 42,19 & 42,38 & 41,71 \\
2025-09 & 582,61 & 591,60 & 581,47 & 42,41 & 42,86 & 42,36 \\
2025-10 & 593,34 & 615,05 & 593,09 & 43,96 & 45,70 & 44,16 \\
2025-11 & 597,86 & 600,90 & 595,25 & 47,87 & 47,46 & 47,78 \\
2025-12 & 577,66 & 593,86 & 574,66 & 45,10 & 45,93 & 45,16 \\
\end{longtable}

Bias-analysen av den tunede modellen viser et systematisk sesongavvik. August og september er underestimerte med
henholdsvis \(-5{,}22\%\) (\(-20\,548\) salgsenheter) og \(-2{,}77\%\) (\(-16\,334\) salgsenheter), mens november og
desember er lett overestimerte med \(+1{,}05\%\) og \(+2{,}36\%\). En mulig forklaring er at kalenderfeaturene ikke
fullt ut fanger en høstoppgang som skiller seg fra tidligere år -- antakelsen er at treningsperioden (2022--2024)
inneholder en annen august--september-profil enn 2025, og at modellen dermed glatter ut det virkelige nivået i disse
månedene. Mønsteret er konsistent på tvers av alle tre modellsporene og peker på en begrensning i det tilgjengelige
sesongsignalet, ikke bare på den tunede modellen.

Figur 7.1 illustrerer den månedlige biasen per modell.

\begin{figure}
\centering
\includegraphics[width=0.8\textwidth,height=\textheight]{../006 analysis/aktiviteter/16_tolke_modellresultater/fig_bias_maaned_modell.png}
\caption{Figur 7.1 Månedlig bias per modell i 2025. Positive verdier betyr at modellen overestimerer salget, negative at
den underestimerer. Kurven viser at alle tre modellene underestimerer i august--september.}
\end{figure}

Figur 7.2 viser samme månedsforløp som Tabell 7.1, men med RMSE visualisert som linjediagram.

\begin{figure}
\centering
\includegraphics[width=0.8\textwidth,height=\textheight]{../006 analysis/aktiviteter/13_rmse_og_mape/fig_rmse_maaned_modell.png}
\caption{Figur 7.2 Månedlig RMSE per modell i 2025. Tuned Random Forest har lavest RMSE i majoriteten av månedene, mens
benchmark lineær slår gjennom i enkeltmåneder.}
\end{figure}

Segmentanalysen (jf. Tabell 7.2) forsterker forskjellen mellom de to metrikkene. Tuned Random Forest vinner RMSE i 13 av
14 tolkningssegmenter -- alle fire kvartaler, alle tre rabattband, alle fire regioner og to av tre salgsnivå -- men må
vike for benchmark lineær og baseline Random Forest i flere av MAPE-segmentene, særlig i høyrabattsegmentet,
vest-regionen og lavvolumssegmentet. I segmentet «høyt salg» er benchmark lineær best på begge metrikker (RMSE
\(699{,}10\), MAPE \(30{,}31\%\)), mens tuned Random Forest likevel gir lavest RMSE i segmentet «lavt salg» til tross
for at tuned Random Forest-MAPE i samme segment ligger på \(90{,}27\,\%\).

Tabell 7.2 gir detaljert RMSE og MAPE per modell for alle 14 tolkningssegmentene.

\begin{longtable}[]{@{}
  >{\raggedright\arraybackslash}p{(\columnwidth - 14\tabcolsep) * \real{0.1250}}
  >{\raggedright\arraybackslash}p{(\columnwidth - 14\tabcolsep) * \real{0.1250}}
  >{\raggedright\arraybackslash}p{(\columnwidth - 14\tabcolsep) * \real{0.1250}}
  >{\raggedright\arraybackslash}p{(\columnwidth - 14\tabcolsep) * \real{0.1250}}
  >{\raggedright\arraybackslash}p{(\columnwidth - 14\tabcolsep) * \real{0.1250}}
  >{\raggedright\arraybackslash}p{(\columnwidth - 14\tabcolsep) * \real{0.1250}}
  >{\raggedright\arraybackslash}p{(\columnwidth - 14\tabcolsep) * \real{0.1250}}
  >{\raggedright\arraybackslash}p{(\columnwidth - 14\tabcolsep) * \real{0.1250}}@{}}
\caption{Tabell 7.2 RMSE og MAPE per modell i utvalgte segmenter (kvartal, rabattbånd, region,
salgsbånd).}\tabularnewline
\toprule\noalign{}
\begin{minipage}[b]{\linewidth}\raggedright
Dimensjon
\end{minipage} & \begin{minipage}[b]{\linewidth}\raggedright
Verdi
\end{minipage} & \begin{minipage}[b]{\linewidth}\raggedright
RMSE lineær
\end{minipage} & \begin{minipage}[b]{\linewidth}\raggedright
RMSE baseline RF
\end{minipage} & \begin{minipage}[b]{\linewidth}\raggedright
RMSE tuned RF
\end{minipage} & \begin{minipage}[b]{\linewidth}\raggedright
MAPE lineær (\%)
\end{minipage} & \begin{minipage}[b]{\linewidth}\raggedright
MAPE baseline RF (\%)
\end{minipage} & \begin{minipage}[b]{\linewidth}\raggedright
MAPE tuned RF (\%)
\end{minipage} \\
\midrule\noalign{}
\endfirsthead
\toprule\noalign{}
\begin{minipage}[b]{\linewidth}\raggedright
Dimensjon
\end{minipage} & \begin{minipage}[b]{\linewidth}\raggedright
Verdi
\end{minipage} & \begin{minipage}[b]{\linewidth}\raggedright
RMSE lineær
\end{minipage} & \begin{minipage}[b]{\linewidth}\raggedright
RMSE baseline RF
\end{minipage} & \begin{minipage}[b]{\linewidth}\raggedright
RMSE tuned RF
\end{minipage} & \begin{minipage}[b]{\linewidth}\raggedright
MAPE lineær (\%)
\end{minipage} & \begin{minipage}[b]{\linewidth}\raggedright
MAPE baseline RF (\%)
\end{minipage} & \begin{minipage}[b]{\linewidth}\raggedright
MAPE tuned RF (\%)
\end{minipage} \\
\midrule\noalign{}
\endhead
\bottomrule\noalign{}
\endlastfoot
Kvartal & 1 & 570,43 & 580,73 & 569,10 & 43,58 & 43,73 & 43,67 \\
Kvartal & 2 & 578,10 & 576,41 & 573,05 & 44,69 & 43,11 & 43,90 \\
Kvartal & 3 & 580,18 & 591,95 & 579,59 & 42,48 & 42,67 & 42,17 \\
Kvartal & 4 & 589,94 & 602,59 & 587,87 & 45,85 & 46,46 & 45,89 \\
Rabattbånd & lav & 583,01 & 598,95 & 581,12 & 46,66 & 46,93 & 46,14 \\
Rabattbånd & middels & 577,89 & 587,10 & 577,11 & 44,51 & 44,63 & 44,37 \\
Rabattbånd & høy & 581,41 & 586,74 & 577,96 & 42,76 & 42,36 & 42,62 \\
Region & Central & 585,83 & 595,13 & 584,92 & 44,36 & 44,49 & 44,39 \\
Region & East & 579,39 & 587,92 & 575,30 & 44,19 & 44,26 & 43,86 \\
Region & South & 587,59 & 599,43 & 587,02 & 45,48 & 45,29 & 45,03 \\
Region & West & 573,90 & 581,36 & 571,78 & 43,44 & 43,18 & 43,27 \\
Salgsbånd & lavt salg & 695,06 & 694,19 & 690,44 & 90,79 & 89,39 & 90,27 \\
Salgsbånd & middels salg & 196,79 & 224,70 & 192,51 & 11,40 & 12,30 & 11,20 \\
Salgsbånd & høyt salg & 699,10 & 713,73 & 699,57 & 30,31 & 30,60 & 30,39 \\
\end{longtable}

Feature importance (jf. Tabell 8.3) støtter tolkningen av hvilke dimensjoner som driver variasjonen: kalendervariablene
utgjør samlet \(\sim 42\%\) av importance i topp 10, \splitcode{Discount} er sterkeste enkeltvariabel med \(11{,}48\%\),
og regionvariablene bidrar med \(\sim 5{,}8\%\). Den lineære modellen gir negativt fortegn på \splitcode{Discount}
(koeffisient \(-166{,}35\)), mens Random Forest rangerer samme variabel øverst. Forskjellen tolkes som et mulig tegn på
en ikke-lineær eller interaktiv sammenheng mellom rabatt og salg som en lineær spesifikasjon ikke klarer å fange opp.

\begin{center}\rule{0.5\linewidth}{0.5pt}\end{center}

\hypertarget{resultat}{%
\subsection{8 Resultat}\label{resultat}}

Tabell 8.1 oppsummerer den samlede ytelsen på 2025 for de tre modellsporene. Tuned Random Forest har lavest RMSE og
MAPE, mens baseline Random Forest er svakest på RMSE og benchmark lineær er svakest på MAPE.

\begin{longtable}[]{@{}
  >{\raggedright\arraybackslash}p{(\columnwidth - 4\tabcolsep) * \real{0.3333}}
  >{\raggedright\arraybackslash}p{(\columnwidth - 4\tabcolsep) * \real{0.3333}}
  >{\raggedright\arraybackslash}p{(\columnwidth - 4\tabcolsep) * \real{0.3333}}@{}}
\caption{Tabell 8.1 Samlet prognoseytelse på 2025 (3 312 observasjoner).}\tabularnewline
\toprule\noalign{}
\begin{minipage}[b]{\linewidth}\raggedright
Modell
\end{minipage} & \begin{minipage}[b]{\linewidth}\raggedright
RMSE
\end{minipage} & \begin{minipage}[b]{\linewidth}\raggedright
MAPE
\end{minipage} \\
\midrule\noalign{}
\endfirsthead
\toprule\noalign{}
\begin{minipage}[b]{\linewidth}\raggedright
Modell
\end{minipage} & \begin{minipage}[b]{\linewidth}\raggedright
RMSE
\end{minipage} & \begin{minipage}[b]{\linewidth}\raggedright
MAPE
\end{minipage} \\
\midrule\noalign{}
\endhead
\bottomrule\noalign{}
\endlastfoot
Benchmark lineær & 580,39 & 44,18 \% \\
Baseline Random Forest & 589,28 & 44,12 \% \\
Tuned Random Forest & 578,26 & 43,97 \% \\
\end{longtable}

Figur 8.1 visualiserer den samlede ytelsen per modell.

\begin{figure}
\centering
\includegraphics[width=0.8\textwidth,height=\textheight]{../006 analysis/aktiviteter/13_rmse_og_mape/fig_rmse_mape_samlet.png}
\caption{Figur 8.1 Samlet RMSE og MAPE for benchmark lineær, Random Forest baseline og tuned Random Forest i 2025. Tuned
Random Forest har lavest RMSE, mens MAPE er tettere fordelt mellom modellene.}
\end{figure}

Tabell 8.2 viser hvor mange av de tolv månedene i 2025 hver modell vinner på henholdsvis RMSE og MAPE.

\begin{longtable}[]{@{}
  >{\raggedright\arraybackslash}p{(\columnwidth - 4\tabcolsep) * \real{0.3333}}
  >{\raggedright\arraybackslash}p{(\columnwidth - 4\tabcolsep) * \real{0.3333}}
  >{\raggedright\arraybackslash}p{(\columnwidth - 4\tabcolsep) * \real{0.3333}}@{}}
\caption{Tabell 8.2 Månedlig vinnertelling per metrikk i 2025.}\tabularnewline
\toprule\noalign{}
\begin{minipage}[b]{\linewidth}\raggedright
Modell
\end{minipage} & \begin{minipage}[b]{\linewidth}\raggedright
RMSE-vinner (av 12)
\end{minipage} & \begin{minipage}[b]{\linewidth}\raggedright
MAPE-vinner (av 12)
\end{minipage} \\
\midrule\noalign{}
\endfirsthead
\toprule\noalign{}
\begin{minipage}[b]{\linewidth}\raggedright
Modell
\end{minipage} & \begin{minipage}[b]{\linewidth}\raggedright
RMSE-vinner (av 12)
\end{minipage} & \begin{minipage}[b]{\linewidth}\raggedright
MAPE-vinner (av 12)
\end{minipage} \\
\midrule\noalign{}
\endhead
\bottomrule\noalign{}
\endlastfoot
Benchmark lineær & 1 & 3 \\
Baseline Random Forest & 0 & 6 \\
Tuned Random Forest & 11 & 3 \\
\end{longtable}

Tabell 8.3 viser de ti variablene med høyest feature importance i den tunede Random Forest-modellen, gruppert etter
variabelgruppe.

\begin{longtable}[]{@{}
  >{\raggedright\arraybackslash}p{(\columnwidth - 6\tabcolsep) * \real{0.2500}}
  >{\raggedright\arraybackslash}p{(\columnwidth - 6\tabcolsep) * \real{0.2500}}
  >{\raggedright\arraybackslash}p{(\columnwidth - 6\tabcolsep) * \real{0.2500}}
  >{\raggedright\arraybackslash}p{(\columnwidth - 6\tabcolsep) * \real{0.2500}}@{}}
\caption{Tabell 8.3 Topp 10 feature importance for tuned Random Forest.}\tabularnewline
\toprule\noalign{}
\begin{minipage}[b]{\linewidth}\raggedright
Rang
\end{minipage} & \begin{minipage}[b]{\linewidth}\raggedright
Variabel
\end{minipage} & \begin{minipage}[b]{\linewidth}\raggedright
Gruppe
\end{minipage} & \begin{minipage}[b]{\linewidth}\raggedright
Importance
\end{minipage} \\
\midrule\noalign{}
\endfirsthead
\toprule\noalign{}
\begin{minipage}[b]{\linewidth}\raggedright
Rang
\end{minipage} & \begin{minipage}[b]{\linewidth}\raggedright
Variabel
\end{minipage} & \begin{minipage}[b]{\linewidth}\raggedright
Gruppe
\end{minipage} & \begin{minipage}[b]{\linewidth}\raggedright
Importance
\end{minipage} \\
\midrule\noalign{}
\endhead
\bottomrule\noalign{}
\endlastfoot
1 & Discount & Pris/kampanje & 11,48 \% \\
2 & dayofmonth & Kalender & 11,35 \% \\
3 & weekofyear & Kalender & 10,40 \% \\
4 & month & Kalender & 6,74 \% \\
5 & dayofweek & Kalender & 6,52 \% \\
6 & year & Kalender & 4,02 \% \\
7 & quarter & Kalender & 2,95 \% \\
8 & Region\_East & Region & 2,04 \% \\
9 & Region\_West & Region & 1,91 \% \\
10 & Region\_Central & Region & 1,84 \% \\
\end{longtable}

Figur 8.2 visualiserer topp 10-rangeringen som horisontalt søyleplot.

\begin{figure}
\centering
\includegraphics[width=0.8\textwidth,height=\textheight]{../006 analysis/aktiviteter/15_viktige_variabler/fig_feature_importance_tuned_top10.png}
\caption{Figur 8.2 Topp 10 feature importance for tuned Random Forest. Variablene er rangert etter normalisert viktighet
og viser hvilke signaler modellen faktisk vektlegger.}
\end{figure}

Figur 8.3 aggregerer samme topp 10 etter variabelgruppe og viser hvilke typer signaler som dominerer prediksjonen.

\begin{figure}
\centering
\includegraphics[width=0.8\textwidth,height=\textheight]{../006 analysis/aktiviteter/15_viktige_variabler/fig_variabelgrupper_tuned_top10.png}
\caption{Figur 8.3 Samlet feature importance per variabelgruppe i topp 10 for tuned Random Forest. Figuren oppsummerer
hvilke grupper av signaler (kalender, pris/kampanje, region) som dominerer modellens prediksjon.}
\end{figure}

Tabell 8.4 oppsummerer vinnermodellen per segment på RMSE og MAPE, fordelt på kvartal, rabattband, region og salgsnivå.

\begin{longtable}[]{@{}
  >{\raggedright\arraybackslash}p{(\columnwidth - 6\tabcolsep) * \real{0.2500}}
  >{\raggedright\arraybackslash}p{(\columnwidth - 6\tabcolsep) * \real{0.2500}}
  >{\raggedright\arraybackslash}p{(\columnwidth - 6\tabcolsep) * \real{0.2500}}
  >{\raggedright\arraybackslash}p{(\columnwidth - 6\tabcolsep) * \real{0.2500}}@{}}
\caption{Tabell 8.4 Vinnermodell per segment på RMSE og MAPE.}\tabularnewline
\toprule\noalign{}
\begin{minipage}[b]{\linewidth}\raggedright
\splitword{Segmentdimensjon}
\end{minipage} & \begin{minipage}[b]{\linewidth}\raggedright
Verdi
\end{minipage} & \begin{minipage}[b]{\linewidth}\raggedright
RMSE-vinner
\end{minipage} & \begin{minipage}[b]{\linewidth}\raggedright
MAPE-vinner
\end{minipage} \\
\midrule\noalign{}
\endfirsthead
\toprule\noalign{}
\begin{minipage}[b]{\linewidth}\raggedright
\splitword{Segmentdimensjon}
\end{minipage} & \begin{minipage}[b]{\linewidth}\raggedright
Verdi
\end{minipage} & \begin{minipage}[b]{\linewidth}\raggedright
RMSE-vinner
\end{minipage} & \begin{minipage}[b]{\linewidth}\raggedright
MAPE-vinner
\end{minipage} \\
\midrule\noalign{}
\endhead
\bottomrule\noalign{}
\endlastfoot
Kvartal & Q1 & Tuned RF & Benchmark lineær \\
Kvartal & Q2 & Tuned RF & Baseline RF \\
Kvartal & Q3 & Tuned RF & Tuned RF \\
Kvartal & Q4 & Tuned RF & Benchmark lineær \\
Rabatt & Lav & Tuned RF & Tuned RF \\
Rabatt & Middels & Tuned RF & Tuned RF \\
Rabatt & Høy & Tuned RF & Baseline RF \\
Region & Central & Tuned RF & Benchmark lineær \\
Region & East & Tuned RF & Tuned RF \\
Region & South & Tuned RF & Tuned RF \\
Region & West & Tuned RF & Baseline RF \\
Salgsnivå & Lavt salg & Tuned RF & Baseline RF \\
Salgsnivå & Middels salg & Tuned RF & Tuned RF \\
Salgsnivå & Høyt salg & Benchmark lineær & Benchmark lineær \\
\end{longtable}

\begin{center}\rule{0.5\linewidth}{0.5pt}\end{center}

\hypertarget{diskusjon}{%
\subsection{9 Diskusjon}\label{diskusjon}}

\hypertarget{tolkning-av-hovedfunn-mot-problemstillingen}{%
\subsubsection{9.1 Tolkning av hovedfunn mot problemstillingen}\label{tolkning-av-hovedfunn-mot-problemstillingen}}

Problemstillingen spør hvordan multippel lineær regresjon og Random Forest Regressor kan brukes til å predikere salg i
2025, og hvilke faktorer som påvirker salget mest. Resultatene gir et differensiert svar. Tuned Random Forest er samlet
best i 2025 med RMSE \(578{,}26\) og MAPE \(43{,}97\%\), men gapet til benchmark lineær er marginalt på totalnivå --
\(2{,}13\) enheter lavere RMSE og \(0{,}21\) prosentpoeng lavere MAPE. Den samlede ytelsen underslår dermed et
tydeligere månedsbilde: tuned Random Forest vinner RMSE i elleve av tolv måneder, men bare tre måneder på MAPE, mens
baseline Random Forest vinner MAPE i seks måneder uten å vinne RMSE en eneste måned. Det betyr at «best samlet» i
praksis handler om absolutt presisjon i måneder med store salgsvolumer, mens prosentfeilen fortsatt er heterogen på
tvers av året.

Bias-mønsteret nyanserer dette ytterligere. Tuned Random Forest underestimerer systematisk i august og september
(henholdsvis \(-5{,}22\%\) og \(-2{,}77\%\)) og overestimerer svakt i november og desember (\(+1{,}05\%\) og
\(+2{,}36\%\)). Mønsteret er konsistent for alle tre modellsporene, noe som peker på en begrensning i sesongsignalet i
datagrunnlaget snarere enn i den enkelte modellen. Samtidig viser segmentanalysen at benchmark lineær er bedre enn tuned
Random Forest i segmentet for høyt salgsnivå på både RMSE og MAPE. Det første delproblemet bør derfor besvares med en
differensiert anbefaling: tuned Random Forest er mest hensiktsmessig som hovedmodell når absolutt prognosepresisjon
prioriteres, mens benchmark lineær bør brukes som supplement i toppbelastningssituasjoner og som forklaringsstøtte.

Tabell 9.1 sammenstiller modellprofilene og tydeliggjør differensieringen.

\begin{landscape}

\begin{longtable}[]{@{}
  >{\raggedright\arraybackslash}p{(\columnwidth - 18\tabcolsep) * \real{0.1000}}
  >{\raggedright\arraybackslash}p{(\columnwidth - 18\tabcolsep) * \real{0.1000}}
  >{\raggedright\arraybackslash}p{(\columnwidth - 18\tabcolsep) * \real{0.1000}}
  >{\raggedright\arraybackslash}p{(\columnwidth - 18\tabcolsep) * \real{0.1000}}
  >{\raggedright\arraybackslash}p{(\columnwidth - 18\tabcolsep) * \real{0.1000}}
  >{\raggedright\arraybackslash}p{(\columnwidth - 18\tabcolsep) * \real{0.1000}}
  >{\raggedright\arraybackslash}p{(\columnwidth - 18\tabcolsep) * \real{0.1000}}
  >{\raggedright\arraybackslash}p{(\columnwidth - 18\tabcolsep) * \real{0.1000}}
  >{\raggedright\arraybackslash}p{(\columnwidth - 18\tabcolsep) * \real{0.1000}}
  >{\raggedright\arraybackslash}p{(\columnwidth - 18\tabcolsep) * \real{0.1000}}@{}}
\caption{Tabell 9.1 Modellprofil med samlet RMSE/MAPE, antall vinnermåneder og -segmenter, tolkbarhetsnivå samt
hovedstyrke og hovedsvakhet per modell.}\tabularnewline
\toprule\noalign{}
\begin{minipage}[b]{\linewidth}\raggedright
\splitword{Modellrolle}
\end{minipage} & \begin{minipage}[b]{\linewidth}\raggedright
RMSE 2025
\end{minipage} & \begin{minipage}[b]{\linewidth}\raggedright
MAPE 2025 (\%)
\end{minipage} & \begin{minipage}[b]{\linewidth}\raggedright
RMSE-vinnermåneder
\end{minipage} & \begin{minipage}[b]{\linewidth}\raggedright
MAPE-vinnermåneder
\end{minipage} & \begin{minipage}[b]{\linewidth}\raggedright
RMSE-vinnersegmenter
\end{minipage} & \begin{minipage}[b]{\linewidth}\raggedright
MAPE-vinnersegmenter
\end{minipage} & \begin{minipage}[b]{\linewidth}\raggedright
Tolkbarhet
\end{minipage} & \begin{minipage}[b]{\linewidth}\raggedright
\splitword{Hovedstyrke}
\end{minipage} & \begin{minipage}[b]{\linewidth}\raggedright
\splitword{Hovedsvakhet}
\end{minipage} \\
\midrule\noalign{}
\endfirsthead
\toprule\noalign{}
\begin{minipage}[b]{\linewidth}\raggedright
\splitword{Modellrolle}
\end{minipage} & \begin{minipage}[b]{\linewidth}\raggedright
RMSE 2025
\end{minipage} & \begin{minipage}[b]{\linewidth}\raggedright
MAPE 2025 (\%)
\end{minipage} & \begin{minipage}[b]{\linewidth}\raggedright
RMSE-vinnermåneder
\end{minipage} & \begin{minipage}[b]{\linewidth}\raggedright
MAPE-vinnermåneder
\end{minipage} & \begin{minipage}[b]{\linewidth}\raggedright
RMSE-vinnersegmenter
\end{minipage} & \begin{minipage}[b]{\linewidth}\raggedright
MAPE-vinnersegmenter
\end{minipage} & \begin{minipage}[b]{\linewidth}\raggedright
Tolkbarhet
\end{minipage} & \begin{minipage}[b]{\linewidth}\raggedright
\splitword{Hovedstyrke}
\end{minipage} & \begin{minipage}[b]{\linewidth}\raggedright
\splitword{Hovedsvakhet}
\end{minipage} \\
\midrule\noalign{}
\endhead
\bottomrule\noalign{}
\endlastfoot
benchmark lineær & 580,39 & 44,18 & 1 & 3 & 1 & 4 & høy & Høy tolkbarhet og \splitword{konkurransedyktig} i enkelte
segmenter. & \splitword{Multikollinearitet} og manglende regularisering svekker robust
\splitword{koeffisienttolkning.} \\
baseline RF & 589,28 & 44,12 & 0 & 6 & 0 & 4 & middels & Sterk lokal MAPE-ytelse og nyttig RF-sammenligningspunkt. &
Svakest samlet og uten RMSE-seire i måneder eller segmenter. \\
tuned RF & 578,26 & 43,97 & 11 & 3 & 13 & 6 & middels & Best samlet 2025 og sterkest på absolutt feil. & MAPE er mer
ujevn enn RMSE mellom måneder og segmenter. \\
\end{longtable}

\end{landscape}

\hypertarget{variablenes-puxe5virkning}{%
\subsubsection{9.2 Variablenes påvirkning}\label{variablenes-puxe5virkning}}

Det andre delproblemet gjelder hvilke faktorer som påvirker salget mest. Rangeringen av feature importance i tuned
Random Forest viser at rabatt og kalendervariabler dominerer bildet. \splitcode{Discount} er det sterkeste
enkeltsignalet med \(11{,}48\%\), tett fulgt av kalendervariablene \splitcode{dayofmonth} (\(11{,}35\%\)),
\splitcode{weekofyear} (\(10{,}40\%\)), \splitcode{month} (\(6{,}74\%\)) og \splitcode{dayofweek} (\(6{,}52\%\)).
Kalendergruppen står samlet for om lag \(42\%\) av importance i topp 10, mens regionvariablene bidrar med \(5{,}8\%\).
De ni øverste variablene i tuned Random Forest gjenfinnes også i topp 10 for baseline Random Forest, noe som styrker
stabiliteten i rangeringen.

Lineær regresjon peker i samme retning, men med annen fortegning. \splitcode{Discount} (\(-166{,}35\)) har den største
meningsfulle koeffisientverdien; \splitcode{Region\_North}-koeffisienten er isolert sett størst i tallverdi
(\(-277{,}20\)), men bygger på én enkelt observasjon (jf. kap. 5.2) og bør tolkes som usikkerhet snarere enn en
substansiell regioneffekt. Det negative fortegnet for \splitcode{Discount} i den lineære modellen står i kontrast til
dens høye feature importance i Random Forest, der retningen ikke er fiksert. Det mest sannsynlige er at sammenhengen
mellom rabatt og salg er ikke-lineær eller samspillende med andre variabler -- for eksempel at rabatteffekten avhenger
av produktkategori eller sesong -- slik at en lineær spesifikasjon overkompenserer med ett fortegn mens et
tremodellbasert ensemble fanger opp flere betingede mønstre. Forskjellen er derfor ikke motstridende funn, men et tegn
på at modellene svarer ulikt på det samme underliggende signalet. Variabelsignalene er prediktive, ikke kausale, og bør
tolkes som indikatorer på hvilke dimensjoner som bærer prognoseverdien.

\hypertarget{praktisk-nytte-for-dagligvare}{%
\subsubsection{9.3 Praktisk nytte for Dagligvare}\label{praktisk-nytte-for-dagligvare}}

Funnene har fire tydelige praktiske bruksområder i Dagligvare, alle forankret i modellvalget i 9.1 og
variabelrangeringen i 9.2. For innkjøp og overordnet lagerstyring er tuned Random Forest det naturlige standardvalget
fordi modellen er best samlet på absolutt presisjon og gir den mest stabile månedsytelsen. Analysen dokumenterer ikke
lageroptimalisering i snever forstand, men gir et sterkere grunnlag for å treffe bedre på totalnivå i planlagte
bestillinger og dermed redusere risiko for både over- og underbestilling.

Tabell 9.2 oppsummerer beslutningsmatrisen med anbefalt modellrolle, prioritert metrikk og hovedforbehold per
bruksområde.

\begin{longtable}[]{@{}
  >{\raggedright\arraybackslash}p{(\columnwidth - 8\tabcolsep) * \real{0.2000}}
  >{\raggedright\arraybackslash}p{(\columnwidth - 8\tabcolsep) * \real{0.2000}}
  >{\raggedright\arraybackslash}p{(\columnwidth - 8\tabcolsep) * \real{0.2000}}
  >{\raggedright\arraybackslash}p{(\columnwidth - 8\tabcolsep) * \real{0.2000}}
  >{\raggedright\arraybackslash}p{(\columnwidth - 8\tabcolsep) * \real{0.2000}}@{}}
\caption{Tabell 9.2 Beslutningsmatrise og bruksregler: beslutningsområder, anbefalt modellrolle, prioritert metrikk,
praktisk nyttegrad og hovedforbehold.}\tabularnewline
\toprule\noalign{}
\begin{minipage}[b]{\linewidth}\raggedright
\splitword{Beslutningsområde}
\end{minipage} & \begin{minipage}[b]{\linewidth}\raggedright
Anbefalt modell
\end{minipage} & \begin{minipage}[b]{\linewidth}\raggedright
Prioritert metrikk
\end{minipage} & \begin{minipage}[b]{\linewidth}\raggedright
Praktisk nytte
\end{minipage} & \begin{minipage}[b]{\linewidth}\raggedright
\splitword{Hovedforbehold}
\end{minipage} \\
\midrule\noalign{}
\endfirsthead
\toprule\noalign{}
\begin{minipage}[b]{\linewidth}\raggedright
\splitword{Beslutningsområde}
\end{minipage} & \begin{minipage}[b]{\linewidth}\raggedright
Anbefalt modell
\end{minipage} & \begin{minipage}[b]{\linewidth}\raggedright
Prioritert metrikk
\end{minipage} & \begin{minipage}[b]{\linewidth}\raggedright
Praktisk nytte
\end{minipage} & \begin{minipage}[b]{\linewidth}\raggedright
\splitword{Hovedforbehold}
\end{minipage} \\
\midrule\noalign{}
\endhead
\bottomrule\noalign{}
\endlastfoot
Innkjøp og lager & tuned RF & RMSE & høy & Ikke \splitword{lageroptimalisering;} prosentfeil kan være ujevn i
delsegmenter. \\
Kampanje og rabatt & tuned RF, med baseline RF som MAPE-kontroll ved høy rabatt & RMSE med MAPE-kontroll & middels &
\splitword{Rabattsignalene} er prediktive, ikke kausale. \\
Bemanning og \splitword{ressursplanlegging} & tuned RF, med benchmark lineær-kontroll ved svært høyt salgsnivå & RMSE &
middels--høy & Ikke butikk- eller \splitword{skiftoptimalisering;} toppbelastning krever ekstra varsomhet. \\
\splitword{Ledelsesrapportering} & tuned RF som hovedprognose, benchmark lineær som \splitword{forklaringsstøtte} & RMSE
med \splitword{forklaringsstøtte} & middels & Lineær modell skal ikke brukes som kausalt bevis for
\splitword{salgsendringer.} \\
\end{longtable}

For kampanje- og rabattvurderinger anbefales tuned Random Forest som hovedmodell, men med en ekstra kontroll mot MAPE i
segmentet for høy rabatt, der baseline Random Forest er best på prosentfeil. Dette gjenspeiler at \splitcode{Discount}
er det sterkeste enkeltsignalet i modellen, men at prosentavvik er mer følsomt enn absoluttavvik i rabattutsatte
perioder. For aggregert bemannings- og ressursplanlegging gir tuned Random Forest et robust bilde av forventet
aktivitetsnivå per måned og kvartal, men benchmark lineær bør trekkes inn når toppbelastning planlegges, siden den er
sterkere i segmentet for høyt salgsnivå. For ledelsesrapportering fungerer tuned Random Forest som hovedprognose, mens
benchmark lineær brukes som forklaringsstøtte når retningen i sentrale signaler må kommuniseres kortfattet for
ikke-tekniske interessenter. I alle disse bruksområdene må analysens avgrensninger holdes eksplisitt: prognosene er
prediktive, ikke kausale, og bør kombineres med lokal fagkunnskap før de omsettes til operative beslutninger.

\hypertarget{metodiske-begrensninger}{%
\subsubsection{9.4 Metodiske begrensninger}\label{metodiske-begrensninger}}

Flere metodiske forhold setter grenser for hvor langt funnene kan strekkes. Datagrunnlaget er representativt for det
simulerte caset, men er bare ett datasett fra én simulert virksomhet. Generalisering til reelle dagligvarekjeder
forutsetter eget datagrunnlag og ny validering. Analysen inkluderer ikke eksterne makroøkonomiske faktorer som
inflasjon, rente eller konjunkturer, og kan derfor ikke belyse hvordan salget reagerer på slike eksterne sjokk.
Modellomfanget er dessuten avgrenset til multippel lineær regresjon og Random Forest Regressor, slik at andre
metodeklasser -- tidsrekkemodeller, gradient boosting, nevrale nett -- ikke er vurdert.

Tabell 9.3 oppsummerer de metodiske begrensningene med konsekvens for pålitelighet og generaliserbarhet.

\begin{longtable}[]{@{}
  >{\raggedright\arraybackslash}p{(\columnwidth - 6\tabcolsep) * \real{0.2500}}
  >{\raggedright\arraybackslash}p{(\columnwidth - 6\tabcolsep) * \real{0.2500}}
  >{\raggedright\arraybackslash}p{(\columnwidth - 6\tabcolsep) * \real{0.2500}}
  >{\raggedright\arraybackslash}p{(\columnwidth - 6\tabcolsep) * \real{0.2500}}@{}}
\caption{Tabell 9.3 Metodiske begrensninger i studien, med tematisk kategori, beskrivelse og konsekvens for pålitelighet
og generaliserbarhet.}\tabularnewline
\toprule\noalign{}
\begin{minipage}[b]{\linewidth}\raggedright
Tema
\end{minipage} & \begin{minipage}[b]{\linewidth}\raggedright
\splitword{Beskrivelse}
\end{minipage} & \begin{minipage}[b]{\linewidth}\raggedright
Konsekvens for \splitword{pålitelighet}
\end{minipage} & \begin{minipage}[b]{\linewidth}\raggedright
Konsekvens for \splitword{generaliserbarhet}
\end{minipage} \\
\midrule\noalign{}
\endfirsthead
\toprule\noalign{}
\begin{minipage}[b]{\linewidth}\raggedright
Tema
\end{minipage} & \begin{minipage}[b]{\linewidth}\raggedright
\splitword{Beskrivelse}
\end{minipage} & \begin{minipage}[b]{\linewidth}\raggedright
Konsekvens for \splitword{pålitelighet}
\end{minipage} & \begin{minipage}[b]{\linewidth}\raggedright
Konsekvens for \splitword{generaliserbarhet}
\end{minipage} \\
\midrule\noalign{}
\endhead
\bottomrule\noalign{}
\endlastfoot
\splitword{representativitet} & Prosjektet bygger på én simulert virksomhet og ett datasett, ikke flere virksomheter
eller markeder. & Påliteligheten innenfor dette caset kan fortsatt være god, men robustheten mot andre kontekster er
ikke testet. & Funnene kan ikke uten videre overføres til andre bedrifter, regioner eller produktmixer. \\
eksterne faktorer & Eksterne \splitword{makroøkonomiske} forhold som inflasjon, rente og konjunkturer er ikke modellert.
& Prognosene kan være mindre robuste hvis 2025 påvirkes av forhold som ikke finnes i feature-settet. & Resultatene
generaliserer dårligere til perioder der eksterne sjokk spiller en større rolle. \\
modellomfang & Prosjektet sammenligner bare lineær regresjon og Random Forest, ikke andre \splitword{modellfamilier.} &
Valgt modell er best i dette prosjektets kandidatfelt, men ikke nødvendigvis best mulig totalt sett. & Det er begrenset
grunnlag for å generalisere at samme modellfamilie vil være best i andre lignende problemer. \\
\splitword{koeffisienttolkning} & Lineær regresjon er kjørt uten \splitword{regularisering,} og dummykoding kan gi
\splitword{multikollinearitet} i \splitword{koeffisientene.} & Dette svekker påliteligheten i sterke fortolkninger av
\splitword{enkeltkoeffisienter} som om de var stabile effektmål. & \splitword{Koeffisientmønstre} kan endre seg når
datastruktur eller \splitword{kategorifordeling} endres i andre case. \\
\splitword{valideringsvindu} & Tuning av Random Forest bruker kun 2024 som intern \splitword{valideringsperiode.} &
Modellvalget er etterprøvbart, men sensitivitet for andre \splitword{valideringsvinduer} er ikke undersøkt. & Det gir
mindre grunnlag for å generalisere tuningvalgene til andre tidsperioder eller sesongmønstre. \\
kausalitet & Analysen er prediktiv og ikke kausal, slik at viktige variabler ikke kan tolkes som dokumenterte årsaker
til salg. & Det er mer pålitelig å bruke funnene som prognosestøtte enn som bevis for \splitword{årsakssammenhenger.} &
Beslutninger som krever kausal innsikt kan ikke generaliseres direkte fra disse prediktive mønstrene. \\
\end{longtable}

Den lineære modellen brukes med uregularisert OLS på en modellmatrise med 67 features, inkludert mange one-hot-kodede
dummyvariabler. Det gir en betydelig risiko for multikollinearitet, som svekker tolkningsvaliditeten av de enkelte
koeffisientene selv om prediksjonskraften kan være upåvirket. Hyperparametertuningen av Random Forest er basert på ett
valideringsår (2024), slik at valget av tuned konfigurasjon reflekterer mønstrene i ett år framfor en mer robust
kryssvalidering over flere perioder. To planendringer i modellutviklingsfasen er dokumentert i endringsloggen: det
felles treningssteget ble gjort om til en verifisering av treningsgrunnlag og modellsignaler siden begge modellene
allerede var trent i forutgående aktiviteter, og hyperparametertuningen ble avgrenset til Random Forest-sporet alene.
Samlet styrket disse endringene sporbarheten i modellutviklingen, men innebærer at lineær regresjon ikke har gjennomgått
tilsvarende optimalisering.

Funnene er i tillegg prediktive og ikke kausale. Variabelrangeringer og koeffisienter forteller hvilke signaler som er
nyttige for å predikere salg, men ikke hvorfor salget endrer seg. Denne avgrensningen er viktig å holde fast ved når
resultatene oversettes til beslutningsstøtte, for eksempel når rabattsignalet skal tolkes i kampanjearbeidet.

\hypertarget{videre-arbeid}{%
\subsubsection{9.5 Videre arbeid}\label{videre-arbeid}}

Fire typer videre arbeid følger naturlig av disse begrensningene. Et første steg er supplering med fagfellevurderte
kilder innen etterspørselspredikering i dagligvarehandel, som kan styrke både metodevalg og tolkning. Et andre steg er
sensitivitetsanalyse av variabelsignalene -- særlig \splitcode{Discount} og regionvariablene -- slik at tolkningen blir
mer robust mot ikke-lineære sammenhenger. Et tredje steg er bredere validering, enten ved å bruke flere år som
valideringsvinduer eller ved å teste modellene mot reelle salgsdata fra en dagligvarekjede. Et fjerde steg er utvidelse
av modellomfanget til tidsrekke- og gradient-boosting-metoder som kan adressere både det systematiske høstbias-mønsteret
og den heterogene MAPE-ytelsen på en mer presis måte.

\begin{center}\rule{0.5\linewidth}{0.5pt}\end{center}

\hypertarget{konklusjon}{%
\subsection{10 Konklusjon}\label{konklusjon}}

Problemstillingen spør hvordan multippel lineær regresjon og Random Forest Regressor kan brukes til å predikere salg i
2025 for en simulert dagligvarekjede, og hvilke faktorer som påvirker salget mest. På tvers av de tre modellsporene er
tuned Random Forest det samlet beste valget med RMSE \(578{,}26\) og MAPE \(43{,}97\%\) på 2025, og modellen vinner RMSE
i elleve av tolv måneder og tretten av fjorten tolkningssegmenter. Gapet til benchmark lineær er likevel marginalt på
totalnivå, og sistnevnte er best i segmentet for høyt salgsnivå. Anbefalingen er derfor at Dagligvare bruker tuned
Random Forest som standardprognose for innkjøp, lager og aggregert ressursplanlegging, benchmark lineær som
forklaringsstøtte og som supplement i toppbelastningssituasjoner, og baseline Random Forest som kontroll mot prosentfeil
i perioder med høy rabatt.

Rabatt og kalendervariabler er de mest påvirkningsrike prediktorene, der \splitcode{Discount} er det sterkeste
enkeltsignalet og kalendergruppen samlet står for om lag \(42\%\) av importance i topp 10 i den anbefalte modellen.
Regionvariablene bidrar som sekundære prediktorer og driver samtidig variasjon i MAPE mellom markedene. Funnene gjelder
innenfor det simulerte caset og bør tolkes som prediktive signaler, ikke kausale forklaringer; operativ bruk forutsetter
at Dagligvare kombinerer modellprognosene med lokal fagkunnskap og vurderer supplerende validering før de omsettes til
beslutninger i innkjøp, kampanjestyring og ressursplanlegging.

\begin{center}\rule{0.5\linewidth}{0.5pt}\end{center}

\hypertarget{bibliografi}{%
\subsection{11 Bibliografi}\label{bibliografi}}

GeeksforGeeks. (2026a, 6. april). \emph{Linear Regression in Machine Learning}. Hentet 13. april 2026 fra
\url{https://www.geeksforgeeks.org/machine-learning/ml-linear-regression/}

GeeksforGeeks. (2026b, 6. april). \emph{Random Forest Algorithm in Machine Learning}. Hentet 13. april 2026 fra
\url{https://www.geeksforgeeks.org/machine-learning/random-forest-algorithm-in-machine-learning/}

IBM. (u.å.-a). \emph{What is linear regression?}. Hentet 13. april 2026 fra
\url{https://www.ibm.com/think/topics/linear-regression}

IBM. (u.å.-b). \emph{What is random forest?}. Hentet 13. april 2026 fra
\url{https://www.ibm.com/think/topics/random-forest}

\hypertarget{vedlegg}{%
\subsection{12 Vedlegg}\label{vedlegg}}

Dette kapitlet samler referanser til stort eller detaljert datamateriale som ikke er limt direkte inn i rapporten, men
som ligger som kildefiler under \splitcode{006 analysis/}. Vedleggene A1--A7 gir leseren direkte tilgang til de
underliggende analyseartefaktene.

Tabell 12.1 lister vedleggene A1--A7 med innhold og tilhørende kildefil under \splitcode{006 analysis/}.

\begin{longtable}[]{@{}
  >{\raggedright\arraybackslash}p{(\columnwidth - 4\tabcolsep) * \real{0.3333}}
  >{\raggedright\arraybackslash}p{(\columnwidth - 4\tabcolsep) * \real{0.3333}}
  >{\raggedright\arraybackslash}p{(\columnwidth - 4\tabcolsep) * \real{0.3333}}@{}}
\caption{Tabell 12.1 Vedleggsreferanser til stort datamateriale som hører til analysen.}\tabularnewline
\toprule\noalign{}
\begin{minipage}[b]{\linewidth}\raggedright
Vedlegg
\end{minipage} & \begin{minipage}[b]{\linewidth}\raggedright
Innhold
\end{minipage} & \begin{minipage}[b]{\linewidth}\raggedright
Kildefil
\end{minipage} \\
\midrule\noalign{}
\endfirsthead
\toprule\noalign{}
\begin{minipage}[b]{\linewidth}\raggedright
Vedlegg
\end{minipage} & \begin{minipage}[b]{\linewidth}\raggedright
Innhold
\end{minipage} & \begin{minipage}[b]{\linewidth}\raggedright
Kildefil
\end{minipage} \\
\midrule\noalign{}
\endhead
\bottomrule\noalign{}
\endlastfoot
A1 & Radvise prognoser for 2025 for alle tre modeller (3 312 rader) &
\href{../006\%20analysis/aktiviteter/12_prognoser_2025/tab_prognoser_2025_detalj.csv}{\splitword{tab\_prognoser\_2025\_detalj.csv}} \\
A2 & Radvise prognosefeil og absoluttfeil for 2025 &
\href{../006\%20analysis/aktiviteter/13_rmse_og_mape/tab_prognosefeil_2025_detalj.csv}{\splitword{tab\_prognosefeil\_2025\_detalj.csv}} \\
A3 & Full tuning-kandidatgrid for Random Forest &
\href{../006\%20analysis/aktiviteter/11_parameterjustering_random_forest/tab_rf_tuning_kandidater.csv}{\splitword{tab\_rf\_tuning\_kandidater.csv}} \\
A4 & Full feature importance for tuned Random Forest &
\href{../006\%20analysis/aktiviteter/15_viktige_variabler/tab_rf_tuned_feature_importance.csv}{\splitword{tab\_rf\_tuned\_feature\_importance.csv}} \\
A5 & Full \splitword{koeffisienttabell} for lineær regresjon &
\href{../006\%20analysis/aktiviteter/08_lineaer_regresjon/tab_lr_koeffisienter.csv}{\splitword{tab\_lr\_koeffisienter.csv}} \\
A6 & Full RMSE og MAPE per måned i lang form &
\href{../006\%20analysis/aktiviteter/13_rmse_og_mape/tab_rmse_mape_maaned.csv}{\splitword{tab\_rmse\_mape\_maaned.csv}} \\
A7 & Full segmentmetrikk per modell &
\href{../006\%20analysis/aktiviteter/16_tolke_modellresultater/tab_segmentmetrikk_modell.csv}{\splitword{tab\_segmentmetrikk\_modell.csv}} \\
\end{longtable}

\end{document}
